\subsection{Examples}

In this subsection (with the exception of Examples 7 and 8), K denotes a nondiscrete locally compact commutative field; dx denotes a Haar measure on the additive group of K.

Recall that \(\mod x = |x|\) when \(\mathbf{K} = \mathbf{R}\) , \(\mod x = |x|^2\) when \(\mathbf{K} = \mathbf{C}\) , \(\mod x = |x|_p\) when \(\mathbf{K} = \mathbf{Q}_p\) .

Example 1. --- General linear group.

Let A be the algebra \(\mathbf{M}_{n}(\mathbf{K})\) . The group \(A^{*}\) of invertible elements of A is none other than the general linear group \(\mathbf{GL}(n,\mathbf{K})\) . For every \(X\in A\) , the reduced norm \(\mathrm{Nrd}_{\mathrm{A}/\mathrm{K}}(X)\) is det X; consequently \(\mathrm{N}_{\mathrm{A}/\mathrm{K}}(X)=(\det X)^{n}\) (Alg., Ch. VIII, §12, No.~3, Prop. 8; cf.~A, III, §9, No.~3, Example 3). Since \(X\mapsto{}^{t}X\) is an isomorphism of A onto the opposite algebra,

\[
\mathrm{N} _ {\mathrm{A} ^ {0} / \mathrm{K}} (X) = \mathrm{N} _ {\mathrm{A} / \mathrm{K}} (^ {t} X) = \det (^ {t} X) ^ {n} = (\det X) ^ {n}.
\]

Then, Prop. 16 of §1, No.~11 proves that the measure

\[
\operatorname{mod} (\det X) ^ {- n} \cdot \bigotimes_ {i, j} d x _ {i j} \quad (X = (x _ {i j}))\tag{3}
\]

is a left and right Haar measure on \(\mathbf{GL}(n,\mathbf{K})\) .

To determine the relatively invariant measures on \(\mathbf{GL}(n,\mathbf{K})\) , we shall rely on the following lemma:

Lemma 5. --- The continuous representations of \(\mathbf{GL}(n,\mathbf{K})\) in \(C^{*}\) are the mappings of the form \(X \mapsto \chi(\det X)\) , where \(\chi\) is a continuous representation of \(K^{*}\) in \(C^{*}\) .

Such a mapping is obviously a continuous representation of \(\mathbf{GL}(n,\mathbf{K})\) in \(C^{*}\) . Conversely, suppose that \(\psi\) is a continuous representation of \(\mathbf{GL}(n,\mathbf{K})\) in \(C^{*}\) . For \(x \in K^{*}\) , set

\[
\widetilde {x} = \left( \begin{array}{c c c c c} x & & & & \\ & 1 & & 0 & \\ & & 1 & & \\ & 0 & & \ddots & \\ & & & & 1 \end{array} \right)
\]

and \(\chi(x) = \psi(\widetilde{x})\) . Then, for every matrix \(X \in \mathbf{GL}(n, K)\) , we have \((\det X^{-1})^{\sim} \cdot X \in \mathbf{SL}(n, K)\) . Since \(\mathbf{SL}(n, K)\) is the commutator subgroup of \(\mathbf{GL}(n, K)\) (A, III, §8, No.~9, Cor. of Prop. 17), \(\psi((\det X^{-1})^{\sim} \cdot X) = 1\) , whence

\[
\psi (X) = \psi \big ((\det X) ^ {\sim} \big) = \chi (\det X).
\]

This established, Cor. 1 of Prop. 10 of §1, No.~8 proves that the relatively invariant measures on \(\mathbf{GL}(n,\mathbf{K})\) are, up to a constant factor, the measures of the form

\[
\chi (\det X) \cdot \bigotimes_ {i j} d x _ {i j} \quad (X = (x _ {i j})),\tag{4}
\]

where \(\chi\) is a continuous representation of \(\mathbf{K}^*\) in \(\mathbf{C}^*\) .

Example 2. --- Affine group.

For every \(X \in \mathbf{GL}(n, \mathbf{K})\) and every \(x \in K^{n}\) , let \((X, x)\) be the affine linear mapping \(\xi \mapsto X\xi + x\) in \(K^{n}\) . The set of \((X, x)\) is the affine group G of \(K^{n}\) (A, II, §9, No.~4). The set T of translations is a closed normal subgroup of G, canonically isomorphic to \(K^{n}\) ; on the other hand, \(\mathbf{GL}(n, \mathbf{K})\) is a closed subgroup of G, and G is the semi-direct product of \(\mathbf{GL}(n, \mathbf{K})\) and \(T = K^{n}\) . One equips G with the (locally compact) topology for which G is the topological semi-direct product of \(\mathbf{GL}(n, \mathbf{K})\) and T (GT, III, §2, No.~10). One has

\[
(X, x) = (1, x) \cdot (X, 0).
\]

On the other hand, if \(X \in \mathbf{GL}(n, \mathbf{K})\) and \(x \in \mathbf{T}\) then, for every \(\xi \in \mathbf{K}^n\) ,

\[
(X, 0) (1, x) (X, 0) ^ {- 1} \xi = X (X ^ {- 1} \xi + x) = \xi + X x = (1, X x) \xi ,
\]

therefore the automorphism \((1,x)\mapsto(X,0)(1,x)(X,0)^{-1}\) of T has modulus mod(det X) (§1, No.~10, Prop. 15). In view of Example 1 and §2, No.~9, Remark, the measure

\[
\operatorname{mod} (\det X) ^ {- n - 1} \cdot \left(\bigotimes_ {i j} d x _ {i j}\right) \otimes \left(\bigotimes_ {i} d x _ {i}\right) \quad (X = (x _ {i j}), x = (x _ {i}))\tag{5}
\]

is a left Haar measure on G. On the other hand, by Prop. 14 of §2, No.~9,

\[
\Delta_ {\mathrm{G}} ((X, x)) = \Delta_ {\mathbf {G L} (n, K)} (X) \Delta_ {K ^ {n}} (x) (\mathrm{mod} \det X) ^ {- 1},
\]

or

\[
\Delta_ {\mathrm{G}} ((X, x)) = \operatorname{mod} (\det X ^ {- 1}).\tag{6}
\]

Thus, a right Haar measure on G is given by

\[
(\mathrm{mod} \det X) ^ {- n} \cdot \left(\bigotimes_ {i j} d x _ {i j}\right) \otimes \left(\bigotimes_ {i} d x _ {i}\right).\tag{7}
\]

Example 3. --- Strict triangular group.

Let \([1, n]\) be the set of integers \(m\) such that \(1 \leqslant m \leqslant n\) . Let \(J\) be a subset of \([1, n] \times [1, n]\) satisfying the following conditions:

\begin{enumerate}
\def\labelenumi{\arabic{enumi})}
\item
  if \((i,j)\in \mathbf{J}\) then \(i <   j\)
\item
  if \((i,j)\notin J\) then, for every integer \(k\) such that \(i < k < j\) , at least one of the two pairs \((i,k)\) and \((k,j)\) does not belong to \(J\) .
\end{enumerate}

Let \(T_J\) be the set of matrices \(Z = (z_{ij})_{1 \leqslant i \leqslant n, 1 \leqslant j \leqslant n}\) with elements in \(K\) , such that \(z_{ii} = 1\) , and \(z_{ij} = 0\) if \(i \neq j\) and \((i,j) \notin J\) . This is a closed subset of \(\mathbf{GL}(n,K)\) . The mapping \(Z \mapsto (z_{ij})_{(i,j) \in J}\) is a homeomorphism of \(T_J\) onto \(K^s\) (where \(s\) denotes the number of elements of \(J\) ). If \(Z' = (z_{ij}') \in T_J\) , then \(Z'Z = (z_{ij}'')\) with

\[
\begin{array}{c} z _ {i j} ^ {\prime \prime} = z _ {i j} + z _ {i j} ^ {\prime} + \sum_ {i <   h <   j} z _ {i h} ^ {\prime} z _ {h j} \quad \text {for} i <   j, \\ z _ {i j} ^ {\prime \prime} = 0 \text {for} i > j, \quad z _ {i i} ^ {\prime \prime} = 1, \end{array}
\]

whence \(Z^{\prime}Z \in T_{J}\) . If \(T_{J}\) is identified with \(K^{s}\) , then the mapping \(Z \mapsto Z^{\prime}Z\) (for fixed \(Z^{\prime}\) ) is identified with an affine mapping, and its determinant is 1, as one sees by ordering the pairs \((i,j) \in J\) lexicographically and applying the following lemma:

Lemma 6. --- Let L be a totally ordered finite set. For every \(\lambda \in L\) , let \(V_{\lambda}\) be a free module of finite dimension over a commutative ring \(k\) ; for \(\lambda, \mu\) in L such that \(\lambda \leqslant \mu\) , let \(f_{\lambda \mu} \in \operatorname{Hom}_k(V_\mu, V_\lambda)\) . Then the linear mapping

\[
(v _ {\lambda}) _ {\lambda \in \mathrm{L}} \mapsto \left(\sum_ {\mu \geqslant \lambda} f _ {\lambda \mu} (v _ {\mu})\right) _ {\lambda \in \mathrm{L}},
\]

from \(\prod_{\lambda \in L}V_{\lambda}\) into \(\prod_{\lambda \in L}V_{\lambda}\) , has determinant \(\prod_{\lambda \in L}\det f_{\lambda \lambda}\) .

One reduces immediately to the case that L is an interval of integers, and the lemma then follows from A, III, §8, No.~6, formula (31).

If \(Z \in \mathrm{T}_{\mathrm{J}}\) , one then sees that there exists \(Z' \in \mathrm{T}_{\mathrm{J}}\) such that \(Z'Z = I_n\) , whence \(Z' = Z^{-1}\) . Thus, \(\mathrm{T}_{\mathrm{J}}\) is a closed subgroup of \(\mathbf{GL}(n,\mathrm{K})\) . On the other hand, Prop. 15 of §1, No.~10 shows that the measure

\[
\bigotimes_ {(i, j) \in J} d z _ {i j}
\]

is a left Haar measure on \(T_{J}\) . By calculating \(ZZ'\) one sees in the same way that this measure is a right Haar measure on \(T_{J}\) .

There is an analogous result if, in the definition of \(T_{J}\) , the roles of rows and columns are interchanged.

When J is the set of pairs \((i,j)\) such that i\textless j, the group \(T_{J}\) is called the upper strict triangular group of order n over K, and is denoted \(\mathrm{T}_{1}(n,\mathrm{K})\) . Its transpose is called the lower strict triangular group.

Example 4. --- Large triangular group.

Let \(n_1, \ldots, n_r\) be integers \(\geqslant 1\) . Set \(p_k = n_1 + \ldots + n_{k-1}\) and \(n = p_{r+1} = n_1 + \cdots + n_r\) . Let \(I_k\) be the set of integers \(j\) such that \(p_k < j \leqslant p_{k+1}\) , and \(J\) the union of the \(I_k \times I_l\) for \(k < l\) . Let \(G\) be the closed subgroup of \(\mathbf{GL}(n, K)\) whose elements are the matrices \((Z_{kl})_{1 \leqslant k \leqslant r, 1 \leqslant l \leqslant r}\) such that:

\begin{enumerate}
\def\labelenumi{\arabic{enumi})}
\item
  each \(Z_{kl}\) is a matrix \((z_{ij})_{i\in \mathrm{I}_k,j\in \mathrm{I}_l}\) of elements of \(\mathbf{K}\) , with \(n_k\) rows and \(n_l\) columns;
\item
  \(Z_{kl} = 0\) for \(k > l\) ;
\item
  \(Z_{kk} \in \mathbf{GL}(n_k, \mathrm{K})\) for \(1 \leqslant k \leqslant r\) .
\end{enumerate}

The formula for block multiplication

\[
\begin{array}{r} \left( \begin{array}{c c c c} Z _ {1 1} & 0 & \ldots & 0 \\ 0 & Z _ {2 2} & \ldots & 0 \\ \vdots & \vdots & \ddots & \vdots \\ 0 & 0 & \ldots & Z _ {r r} \end{array} \right) \left( \begin{array}{c c c c} 1 & Z _ {1 2} & \ldots & Z _ {1 r} \\ 0 & 1 & \ldots & Z _ {2 r} \\ \vdots & \vdots & \ddots & \vdots \\ 0 & 0 & \ldots & 1 \end{array} \right) \\ = \left( \begin{array}{c c c c} Z _ {1 1} & Z _ {1 1} Z _ {1 2} & \ldots & Z _ {1 1} Z _ {1 r} \\ 0 & Z _ {2 2} & \ldots & Z _ {2 2} Z _ {2 r} \\ \vdots & \vdots & \ddots & \vdots \\ 0 & 0 & \ldots & Z _ {r r} \end{array} \right) \end{array}\tag{8}
\]

shows that G is the topological semi-direct product of the subgroup D of elements \((Z_{kl}) \in G\) such that \(Z_{kl} = 0\) for \(k \neq l\) and the subgroup \(T_{J}\) of Example 3. Moreover, D is isomorphic to the direct product of the groups \(\mathbf{GL}(n_{k}, \mathbf{K})\) for \(1 \leqslant k \leqslant r\) .

Let \(J'\) be the set of pairs \((j, i)\) for \((i, j) \in J\) and let H be the set of pairs \((i, j) \in [1, n] \times [1, n]\) not belonging to \(J'\) . Let \(Z' = (z_{ij})_{1 \leqslant i \leqslant n, 1 \leqslant j \leqslant n}\) be an element of G. By Prop. 14 of §2, No.~9 and the above Examples 1 and 3, one obtains a left Haar measure on G by taking the image of the measure

\[
\bigotimes_ {k = 1} ^ {r} \left(\left(\operatorname{mod} \det Z _ {k k}\right) ^ {- n _ {k}} \cdot \bigotimes_ {i, j \in I _ {k}} d z _ {i j}\right) \otimes \left(\bigotimes_ {(i, j) \in J} d z _ {i j}\right)
\]

under the mapping

\[
\bigl ((Z _ {k k}), (Z _ {k l}) \bigr) \mapsto \left( \begin{array}{c c c c} Z _ {1 1} & Z _ {1 1} Z _ {1 2} & \ldots & Z _ {1 1} Z _ {1 r} \\ 0 & Z _ {2 2} & \ldots & Z _ {2 2} Z _ {2 r} \\ \vdots & \vdots & \ddots & \vdots \\ 0 & 0 & \ldots & Z _ {r r} \end{array} \right).
\]

Now, consider, for \(k < l\) , the vector space of matrices \(Z_{kl} = (z_{ij})_{i \in \mathbf{I}_k, j \in \mathbf{I}_l}\) . It is the direct sum of the \(n_l\) subspaces \(\mathbf{M}_j\) ( \(j \in \mathbf{I}_l\) ) formed by the matrices such that \(z_{ih} = 0\) for \(h \neq j\) . Each of these subspaces \(\mathbf{M}_j\) is stable under the mapping \(Z_{kl} \mapsto Z_{kk}Z_{kl}\) , and the restriction of this mapping to \(\mathbf{M}_j\) has matrix \(Z_{kk}\) . Consequently (§1, No.~10, Prop. 15) the image of the measure

\(\bigotimes_{i \in \mathbf{I}_k, j \in \mathbf{I}_l} dz_{ij}\) under the mapping \(Z_{kl} \mapsto Z_{kk} Z_{kl}\) is

\[
(\mathrm{mod} \det Z _ {k k}) ^ {- n _ {l}} \cdot \bigotimes_ {i \in \mathrm{I} _ {k}, j \in \mathrm{I} _ {l}} d z _ {i j}.
\]

A left Haar measure on G is therefore given by

\[
\prod_ {k = 1} ^ {r} (\mathrm{mod} \det Z _ {k k}) ^ {- q _ {k}} \cdot \bigotimes_ {(i, j) \in \mathrm{H}} d z _ {i j}\tag{9}
\]

with \(q_{k} = \sum_{k\leqslant l\leqslant r}n_{l} = n - p_{k}\)

Let us calculate the modulus of G, again using Prop. 14 of §2. The groups D and T \(_{J}\) are unimodular; on the other hand:

\[
\begin{array}{c} \left( \begin{array}{c c c c} Z _ {1 1} & 0 & \ldots & 0 \\ 0 & Z _ {2 2} & \ldots & 0 \\ \vdots & \vdots & \ddots & \vdots \\ 0 & 0 & \ldots & Z _ {r r} \end{array} \right) \left( \begin{array}{c c c c} 1 & Z _ {1 2} & \ldots & Z _ {1 r} \\ 0 & 1 & \ldots & Z _ {2 r} \\ \vdots & \vdots & \ddots & \vdots \\ 0 & 0 & \ldots & 1 \end{array} \right) \left( \begin{array}{c c c c} Z _ {1 1} & 0 & \ldots & 0 \\ 0 & Z _ {2 2} & \ldots & 0 \\ \vdots & \vdots & \ddots & \vdots \\ 0 & 0 & \ldots & Z _ {r r} \end{array} \right) ^ {- 1} \\ = \left( \begin{array}{c c c c} 1 & Z _ {1 2} ^ {\prime} & \ldots & Z _ {1 r} ^ {\prime} \\ 0 & 1 & \ldots & Z _ {2 r} ^ {\prime} \\ \vdots & \vdots & \ddots & \vdots \\ 0 & 0 & \ldots & 1 \end{array} \right), \end{array}
\]

where \(Z_{kl}^{\prime}=Z_{kk}Z_{kl}Z_{ll}^{-1}\) . Taking into account Example 3, and Prop. 15 of §1, No.~10, and arguing as above, one sees that if \(X=\mathrm{diag}(Z_{11},\ldots,Z_{rr})\in\mathbf{D}\) then the modulus of the automorphism \(Z\mapsto X^{-1}ZX\) of \(T_{J}\) is

\[
\prod_ {k <   l} (\mathrm{mod} \det Z _ {k k}) ^ {- n _ {l}} (\mathrm{mod} \det Z _ {l l}) ^ {n _ {k}},
\]

therefore

\[
\Delta_ {\mathrm{G}} (Z) = \prod_ {k = 1} ^ {r} (\mathrm{mod} \det Z _ {k k}) ^ {n + n _ {k} - 2 q _ {k}}.\tag{10}
\]

The transposed group \(G'\) of G is studied in the same manner. For \(G'\) , one finds as left Haar measure

\[
\prod_ {k = 1} ^ {r} (\mathrm{mod} \det Z _ {k k}) ^ {- p _ {k + 1}} \cdot \bigotimes_ {(j, i) \in \mathrm{H}} d z _ {i j},
\]

and as modulus

\[
\prod_ {k = 1} ^ {r} (\mathrm{mod} \det Z _ {k k}) ^ {n + n _ {k} - 2 p _ {k + 1}}.
\]

If in particular one takes \(n_{1}=\ldots=n_{r}=1\) , one finds as group G the group \(T(n,K)^{*}\) of invertible elements of the subalgebra of \(\mathbf{M}_{n}(\mathbf{K})\) formed by the matrices \(X=(x_{ij})\) such that \(x_{ij}=0\) for i\textgreater j. This algebra, which we shall denote \(\mathrm{T}(n,\mathrm{K})\) , is called the upper triangular algebra, and the group \(\mathrm{T}(n,\mathrm{K})^{*}\) is called the upper large triangular group of order n over K. The preceding formulas then take the following form: a left Haar measure on \(\mathrm{T}(n,\mathrm{K})^{*}\) is

\[
\prod_ {i = 1} ^ {n} (\mathrm{mod} z _ {i i}) ^ {i - n - 1} \cdot \bigotimes_ {i \leqslant j} d z _ {i j} \quad (Z = (z _ {i j}))\tag{9 bis}
\]

and the modulus of \(\mathrm{T}(n,\mathbf{K})^*\) is

(10 bis)

\[
\Delta_ {\mathrm{T} (n, \mathrm{K}) ^ {*}} (Z) = \prod_ {i = 1} ^ {n} (\mathrm{mod} z _ {i i}) ^ {2 i - n - 1} \quad \left(Z = \left(z _ {i j}\right)\right).
\]

For the transpose of \(\mathrm{T}(n,\mathbf{K})^*\) , or lower large triangular group, one finds as left Haar measure

\[
\prod_ {i = 1} ^ {n} (\mathrm{mod} z _ {i i}) ^ {- i} \cdot \bigotimes_ {i \geqslant j} d z _ {i j},
\]

and as modulus

\[
\prod_ {i = 1} ^ {n} (\mathrm{mod} z _ {i i}) ^ {n + 1 - 2 i}.
\]

Remark. --- The group \(\mathrm{T}(n,\mathrm{K})^{*}\) is a closed subgroup of \(\mathbf{GL}(n,\mathrm{K})\) , and \(\Delta_{\mathrm{T}(n,\mathrm{K})^{*}}\left((z_{ij})\right)=\prod_{i=1}^{n}(\mathrm{mod}z_{ii})^{2i-n-1}\) . We saw in Example 1 that \(\Delta_{\mathbf{GL}(n,\mathrm{K})}=1\) . If n\textgreater1, the function

\[
\Delta_ {\mathrm{T} (n, \mathrm{K}) ^ {*}} / \Delta_ {\mathbf {G L} (n, \mathrm{K})}
\]

on \(\mathrm{T}(n,\mathbf{K})^*\) cannot be extended to a continuous representation of \(\mathbf{GL}(n,\mathbf{K})\) in \(\mathbf{C}^*\) (because such a representation would be equal to 1 on \(\mathbf{SL}(n,\mathbf{K})\) by Lemma 5, whereas \(\mathrm{mod}(z_{11})^{1 - n}\neq 1\) for \(z_{11}\) suitably chosen). It follows that the homogeneous space \(\mathbf{GL}(n,\mathbf{K}) / \mathrm{T}(n,\mathbf{K})^*\) admits no relatively invariant measure if \(n > 1\) ( \(\S 2\) , No.~6, Cor. 1 of Th. 3).

This homogeneous space may be identified, for n = 2, with the projective line over K. For, let \((e_{1}, e_{2})\) be the canonical basis of \(K^{2}\) . The group \(\mathbf{GL}(2, \mathbf{K})\) operates transitively on the set of lines of \(K^{2}\) with 0 omitted, and the stabilizer of \(Ke_{1} - \{0\}\) is \(T(2, K)^{*}\) .

Example 5. --- Special triangular group.

Let us take up again the notations at the beginning of Example 4, and consider the subgroup \(\mathbf{G}_{1} = \mathbf{G} \cap \mathbf{SL}(n, \mathbf{K})\) . This subgroup is the topological semi-direct product of the group \(\mathbf{D}_{1} = \mathbf{D} \cap \mathbf{SL}(n, \mathbf{K})\) with \(T_{J}\) . The group \(D_{1}\) has a normal subgroup A isomorphic to \(\mathbf{SL}(n_{r}, \mathbf{K})\) , namely the subgroup consisting of the elements \(\text{diag}(Z_{kk})\) with \(Z_{kk} = 1\) for k \textless{} r. The homomorphism

\[
\varphi : \operatorname{diag} (Z _ {1 1}, \dots , Z _ {r r}) \mapsto (Z _ {1 1}, \dots , Z _ {r - 1, r - 1})
\]

of \(D_{1}\) into \(\mathbf{GL}(n_{1},\mathbf{K})\times\cdots\times\mathbf{GL}(n_{r-1},\mathbf{K})\) is surjective and has kernel A. On the other hand, \(\varphi\) is continuous. Taking into account Lemma 2 of Appendix I, \(D_{1}/A\) may be identified with \(\mathbf{GL}(n_{1},\mathbf{K})\times\cdots\times\mathbf{GL}(n_{r-1},\mathbf{K})\) . We shall denote by \(\mu\) the Haar measure of A (cf.~Example 6) and by

\[
\alpha = \bigotimes_ {k = 1} ^ {r - 1} \left(\left(\operatorname{mod} \det Z _ {k k}\right) ^ {- n _ {k}} \cdot \bigotimes_ {i, j \in I _ {k}} d z _ {i j}\right) \otimes^ {\prime} d \mu \left(Z _ {r r}\right)
\]

the Haar measure on \(\mathbf{D}_1\) such that

\[
\alpha / \mu = \bigotimes_ {k = 1} ^ {r - 1} \left(\left(\operatorname{mod} \det Z _ {k k}\right) ^ {- n _ {k}} \cdot \bigotimes_ {i, j \in I _ {k}} d z _ {i j}\right)
\]

(§2, No.~7, Prop. 10). One then shows as in Example 4 that a left Haar measure on \(\mathbf{G}_1\) is given by

\[
\begin{array}{l} \bmod \Big (\prod_ {k = 1} ^ {r - 1} (\det Z _ {k k}) ^ {n _ {k} - q _ {k}} \Big) \\ \qquad \cdot \left[ \bigotimes_ {k = 1} ^ {r - 1} \big ((\bmod \det Z _ {k k}) ^ {- n _ {k}} \cdot \bigotimes_ {i, j \in I _ {k}} d z _ {i j} \big) \otimes^ {\prime} d \mu (Z _ {r r}) \right] \otimes \bigotimes_ {(i, j) \in J} d z _ {i j}. \end{array}
\]

Since \(\mathbf{G}_1\) is normal in \(\mathbf{G}\) , the modulus of \(\mathbf{G}_1\) is the restriction of that of \(\mathbf{G}\) (§2, No.~7, Prop. 10 b)).

If \(n_{r}=1\) , the subgroup A reduces to the neutral element, and a left Haar measure on G is

\[
\operatorname{mod} \left(\prod_ {k = 1} ^ {r - 1} (\det Z _ {k k}) ^ {- q _ {k}}\right) \cdot \bigotimes_ {k = 1} ^ {r - 1} \left(\bigotimes_ {i, j \in I _ {k}} d z _ {i j}\right) \otimes \bigotimes_ {(i, j) \in J} d z _ {i j}.
\]

If one takes \(n_{1}=n_{2}=\ldots=n_{r}=1\) , the group \(G_{1}\) obtained is called the upper special triangular group and its transpose \(G_{1}^{\prime}\) is called the lower special triangular group. A left Haar measure on \(G_{1}\) is

\[
\operatorname{mod} \left(\prod_ {i = 1} ^ {n - 1} z _ {i i} ^ {i - n - 1}\right) \cdot \left(\bigotimes_ {i = 1} ^ {n - 1} d z _ {i i}\right) \otimes \left(\bigotimes_ {1 \leqslant i <   j \leqslant n} d z _ {i j}\right)\tag{11}
\]

and the modulus of \(\mathbf{G}_1\) is

\[
\operatorname{mod} \left(\prod_ {i = 1} ^ {n - 1} z _ {i i} ^ {2 i - 2 n}\right).\tag{12}
\]

For \(G_{1}^{\prime}\) one finds similarly the left Haar measure

\[
\operatorname{mod} \left(\prod_ {i = 1} ^ {n - 1} z _ {i i} ^ {n - i - 1}\right) \cdot \left(\bigotimes_ {i = 1} ^ {n - 1} d z _ {i i}\right) \otimes \left(\bigotimes_ {1 \leqslant j <   i \leqslant n} d z _ {i j}\right)
\]

and modulus

\[
\operatorname{mod} \biggl (\prod_ {i = 1} ^ {n - 1} z _ {i i} ^ {2 n - 2 i} \biggr).
\]

Example 6. --- Special linear group.

The closed subgroups \(\mathrm{T}_{1}(n,\mathrm{K})\) and \({}^{t}(\mathrm{T}(n,\mathrm{K})^{*})\) of \(\mathbf{GL}(n,\mathbf{K})\) have intersection \(\{e\}\) . Thus the mapping \((M,N)\mapsto M\cdot N\) is a continuous bijection \(\varphi\) of \(\mathrm{T}_{1}(n,\mathrm{K})\times{}^{t}(\mathrm{T}(n,\mathrm{K})^{*})\) onto a subset \(\Omega\) of \(\mathbf{GL}(n,\mathbf{K})\) .

Lemma 7. --- a) Let \(U = (u_{ij}) \in \mathbf{GL}(n, K)\) . In order that \(U \in \Omega\) , it is necessary and sufficient that \(\det(u_{ij})_{k \leqslant i,j \leqslant n} \neq 0\) for \(k = 2, 3, \ldots, n\) .

\begin{enumerate}
\def\labelenumi{\alph{enumi})}
\setcounter{enumi}{1}
\item
  \(\Omega\) is an open subset of \(\mathbf{GL}(n,\mathbf{K})\) .
\item
  The mapping \(\varphi\) is a homeomorphism of \(\mathrm{T}_1(n,\mathrm{K})\times {}^t (\mathrm{T}(n,\mathrm{K})^*)\) onto \(\Omega\) .
\end{enumerate}

In order that \(U \in \Omega\) , it is necessary and sufficient that there exist a \(Z = (z_{ij}) \in \mathrm{T}_1(n,\mathrm{K})\) such that \(ZU \in {}^t (\mathbf{T}(n,\mathbf{K}))\) (then necessarily

\(ZU \in {}^{t}(\mathrm{T}(n,\mathrm{K})^{*})\) since \(U\) and \(Z\) are invertible). By what we saw earlier, if \(Z\) exists then it is unique. Thus, in order that \(U \in \Omega\) , it is necessary and sufficient that the linear system

\[
\sum_ {k = 1} ^ {n} z _ {i k} u _ {k j} = 0 \quad (1 \leqslant i <   j \leqslant n)
\]

(where \((z_{ij})\in \mathrm{T}_1(n,\mathbf{K})\) ) admit a unique solution. Now, this system may be written

\[
\sum_ {k = i + 1} ^ {n} z _ {i k} u _ {k j} = - u _ {i j} \quad (1 \leqslant i <   j \leqslant n).\tag{13}
\]

For fixed \(i\) , one has a system of \(n - i\) equations in the unknowns \(z_{i,i+1}\) , \(z_{i,i+2}\) , \ldots, \(z_{i,n}\) ; for these systems to admit unique solutions, it is necessary and sufficient that

\[
\det (u _ {k j}) _ {i + 1 \leqslant k \leqslant n, i + 1 \leqslant j \leqslant n} \neq 0
\]

for \(i = 1,2,\ldots ,n - 1\) . This proves a). From this it follows that \(\Omega\) is open in \(\mathbf{GL}(n,\mathbf{K})\) . On the other hand, on solving the system (13) by means of Cramer's formulas, the \(z_{ij}\) are obtained as rational functions of the \(u_{ij}\) with nonzero denominators in \(\Omega\) , therefore \(Z\) depends continuously on \(U\) in \(\Omega\) , which proves c).

Now let \(\mathbf{G}_{1}^{\prime}\subset{}^{t}(\mathrm{T}(n,\mathrm{K})^{*})\) be the lower special triangular group. The mapping \((M,N)\mapsto M\cdot N\) is a continuous bijection \(\psi\) of \(\mathrm{T}_{1}(n,\mathrm{K})\times\mathbf{G}_{1}^{\prime}\) onto a subset \(\Omega^{\prime}\) of \(\mathbf{SL}(n,\mathbf{K})\) .

Lemma 8. --- a) Let \(U = (u_{ij}) \in \mathbf{SL}(n, \mathbf{K})\) . In order that \(U \in \Omega'\) , it is necessary and sufficient that \(\det(u_{ij})_{k \leqslant i,j \leqslant n} \neq 0\) for \(k = 2,3,\ldots,n\) .

\begin{enumerate}
\def\labelenumi{\alph{enumi})}
\setcounter{enumi}{1}
\item
  \(\Omega'\) is an open subset of \(\mathbf{SL}(n,\mathbf{K})\) .
\item
  The mapping \(\psi\) is a homeomorphism of \(\mathrm{T}_1(n,\mathbf{K})\times \mathrm{G}_1'\) onto \(\Omega^{\prime}\) . For, let \(M\in \mathrm{T}_1(n,\mathbf{K})\) and \(N\in {}^t (\mathrm{T}(n,\mathbf{K})^*)\) . In order that \(M\cdot N\in \mathbf{SL}(n,\mathbf{K})\) , it is necessary and sufficient that \(N\in \mathbf{G}_1'\) . Therefore \(\Omega^{\prime} = \mathbf{SL}(n,\mathbf{K})\cap \Omega\) and Lemma 8 follows at once from Lemma 7.
\end{enumerate}

PROPOSITION 6. --- a) The group \(\mathbf{SL}(n,\mathbf{K})\) is unimodular.

\begin{enumerate}
\def\labelenumi{\alph{enumi})}
\setcounter{enumi}{1}
\item
  Let \(\mu_1\) and \(\mu_2\) be left Haar measures on the upper strict triangular group \(\mathrm{T}_1(n, \mathrm{K})\) and the lower special triangular group \(\mathrm{G}_1'\) , respectively. The image of \(\mu_1 \otimes \mu_2\) under the homeomorphism \((M, N) \mapsto M \cdot N^{-1}\) of \(\mathrm{T}_1(n, \mathrm{K}) \times \mathrm{G}_1'\) onto \(\Omega'\) is the restriction to \(\Omega'\) of a Haar measure on \(\mathrm{SL}(n, \mathrm{K})\) .
\item
  The complement of \(\Omega'\) in \(\mathbf{SL}(n,\mathbf{K})\) is negligible for the Haar measure of \(\mathbf{SL}(n,\mathbf{K})\) .
\end{enumerate}

The group \(\mathbf{GL}(n,\mathbf{K})\) is unimodular (Example 1), and \(\mathbf{SL}(n,\mathbf{K})\) is a normal subgroup of \(\mathbf{GL}(n,\mathbf{K})\) , hence is unimodular ( \(\S 2\) , No.~7, Prop. 10 b)). The assertion b) follows from a), Lemma 8, and Prop. 13 of \(\S 2\) , No.~9. Let us prove c). By Lemma 8 a), it suffices to prove the following: if \(p((u_{ij})_{1 \leqslant i,j \leqslant n})\) is a polynomial, not identically zero on \(\mathbf{SL}(n,\mathbf{K})\) , then the set E of \(U \in \mathbf{SL}(n,\mathbf{K})\) such that \(p(U) = 0\) is negligible for the Haar measure. Taking into account \(\S 1\) , No.~10, Cor. of Prop. 13, the topology of \(\mathbf{SL}(n,\mathbf{K})\) has a countable base. It therefore suffices to prove that for every \(U_0 \in \mathbf{E}\) , there exists a neighborhood of \(U_0\) in \(\mathbf{SL}(n,\mathbf{K})\) whose intersection with E is negligible; or again that there exists a neighborhood W of I in \(\mathbf{SL}(n,\mathbf{K})\) such that \(U_0^{-1}\mathbf{E} \cap \mathbf{W}\) is negligible. Let us take \(\mathbf{W} = \Omega'\) . In view of b), it all comes down to showing that the set of pairs \((M,N) \in \mathrm{T}_1(n,\mathbf{K}) \times \mathrm{G}_1'\) such that \(p(U_0MN) = 0\) is negligible for \(\mu_1 \otimes \mu_2\) . By the expressions for \(\mu_1\) and \(\mu_2\) (calculated in Examples 3 and 5), this will result from the following lemma:

Lemma 9. --- Let \(\psi\) be a polynomial \(\neq 0\) of \(\mathbf{K}[\mathbf{X}_1, \ldots, \mathbf{X}_r]\) . In the space \(\mathbf{K}^r\) , the set \(\mathbf{N}\) defined by \(\psi(x_1, \ldots, x_r) = 0\) is negligible for the Haar measure.

Let us argue by induction on r. The lemma is evident for r = 1, since N is then a finite set. Changing if necessary the numbering of the variables, we can suppose that \(\psi \notin K[X_{1}, \ldots, X_{r-1}]\) ; write

\[
\psi \left(\mathrm{X} _ {1}, \dots , \mathrm{X} _ {r}\right) = \mathrm{X} _ {r} ^ {m} \psi_ {0} \left(\mathrm{X} _ {1}, \dots , \mathrm{X} _ {r - 1}\right) + \dots + \psi_ {m} \left(\mathrm{X} _ {1}, \dots , \mathrm{X} _ {r - 1}\right)
\]

with \(m > 0\) and \(\psi_0 \neq 0\) . In the space \(\mathbf{K}^{r-1}\) , let \(\mathbf{N}_0\) be the set defined by \(\psi_0(x_1, \ldots, x_{r-1}) = 0\) , which is negligible by the induction hypothesis. For every \((x_1, \ldots, x_{r-1}) \notin \mathbf{N}_0\) , the set of \(x_r \in \mathbf{K}\) such that \((x_1, \ldots, x_{r-1}, x_r) \in \mathbf{N}\) is finite, therefore negligible. Since \(\mathbf{K}^r\) is countable at infinity (§1, No.~10, Cor. of Prop. 13), \(\mathbf{N} \cap [(\mathbf{K}^{r-1} - \mathbf{N}_0) \times \mathbf{K}]\) is negligible in \(\mathbf{K}^r\) (Ch. V, §8, No.~2, Prop. 4). Therefore \(\mathbf{N}\) is negligible.

Example 7. --- Iwasawa decomposition of \(\mathbf{GL}(n,\mathbf{K})\).

In this example, K denotes one of the fields R, C, H. If \(\lambda \in K\) , \(\overline{\lambda}\) is defined to be equal to \(\lambda\) if K = R, and to the conjugate of \(\lambda\) if K = C or H. Let E be a right vector space over K of dimension n, and let \(\Phi\) be a nondegenerate positive hermitian form on E.

Lemma 10. --- Let \((f_1, f_2, \ldots, f_n)\) be a basis of E.

\begin{enumerate}
\def\labelenumi{\alph{enumi})}
\item
  There exists one and only one orthonormal basis \((e_{1}, e_{2}, \ldots, e_{n})\) of E such that \(f_{i} = e_{1}\alpha_{i1} + e_{2}\alpha_{i2} + \cdots + e_{i}\alpha_{ii}\) ( \(i = 1, 2, \ldots, n\) ) with \(\alpha_{ii} > 0\) for all i.
\item
  For fixed \(\Phi\) , the \(e_i\) and \(\alpha_{ij}\) depend continuously on \((f_1, \ldots, f_n) \in \mathbf{E}^n\) .
\end{enumerate}

Let \(E_{i}=f_{1}K+f_{2}K+\cdots+f_{i}K\) , which has dimension i. Let \(g_{i}\) be a nonzero element of \(E_{i}\) orthogonal to \(E_{i-1}\) and such that \(\Phi(g_{i},g_{i})=1\) . By induction on i, one sees that \((g_{1},\ldots,g_{i})\) is an orthonormal basis of \(E_{i}\) . In particular, \((g_{1},\ldots,g_{n})\) is an orthonormal basis of E. Let \(\lambda_{i}=\Phi(f_{i},g_{i})\) . Since \(f_{i}\notin E_{i-1}\) , one has \(\lambda_{i}\neq0\) . Set \(e_{i}=g_{i}\left|\lambda_{i}\right|\lambda_{i}^{-1}\) . Then

\[
\Phi (e _ {i}, e _ {i}) = | \lambda_ {i} | ^ {2} \overline {{{{\lambda}}}} _ {i} ^ {- 1} \Phi (g _ {i}, g _ {i}) \lambda_ {i} ^ {- 1} = 1,
\]

thus \((e_1, \ldots, e_i)\) is also an orthonormal basis of \(\mathbf{E}_i\) ; moreover, \(\Phi(e_i, f_i) = |\lambda_i| \overline{\lambda}_i^{-1} \Phi(g_i, f_i) = |\lambda_i| > 0\) , thus the \(e_i\) have the properties of a). Let \((e_1', \ldots, e_n')\) be another orthonormal basis of \(\mathbf{E}\) with the same properties. One sees by induction on \(i\) that \((e_1', \ldots, e_i')\) must be a basis of \(\mathbf{E}_i\) , therefore \(e_i' = e_i \mu_i\) for some \(\mu_i \in \mathbf{K}\) . Then

\[
1 = \Phi (e _ {i} ^ {\prime}, e _ {i} ^ {\prime}) = \overline {{\mu}} _ {i} \Phi (e _ {i}, e _ {i}) \mu_ {i} = \overline {{\mu}} _ {i} \mu_ {i},
\]

and \(0 < \Phi(e_i', f_i) = \overline{\mu}_i \Phi(e_i, f_i)\) , therefore \(\mu_i > 0\) and \(\mu_i^2 = 1\) , thus \(\mu_i = 1\) , whence a). Suppose already proven that the \(e_i\) and \(\alpha_{ij}\) depend continuously on \((f_1, \ldots, f_n)\) for \(i < i_0\) , and let us prove that \(e_{i_0}\) and the \(\alpha_{i_0j}\) depend continuously on \((f_1, \ldots, f_n)\) . For \(j < i_0\) , \(\overline{\alpha}_{i_0j} = \Phi(f_{i_0}, e_j)\) depends continuously on \((f_1, \ldots, f_n)\) by the induction hypothesis. On the other hand,

\[
\Phi (f _ {i _ {0}}, f _ {i _ {0}}) = | \alpha_ {i _ {0} 1} | ^ {2} + | \alpha_ {i _ {0} 2} | ^ {2} + \dots + | \alpha_ {i _ {0}, i _ {0} - 1} | ^ {2} + \alpha_ {i _ {0} i _ {0}} ^ {2},
\]

thus \(\alpha_{i_0i_0}\) depends continuously on \((f_1,\ldots ,f_n)\) . Therefore

\[
e _ {i _ {0}} = (f _ {i _ {0}} - e _ {1} \alpha_ {i _ {0} 1} - \dots - e _ {i _ {0} - 1} \alpha_ {i _ {0}, i _ {0} - 1}) \alpha_ {i _ {0} i _ {0}} ^ {- 1}
\]

depends continuously on \((f_1, \ldots, f_n)\) .

Henceforth let \(\mathbf{E} = \mathbf{K}^n\) and let us take for \(\Phi\) the form

\[
\overline {{{x}}} _ {1} y _ {1} + \dots + \overline {{{x}}} _ {n} y _ {n}.
\]

Recall that \(\mathbf{U}(n,\mathbf{K})\) denotes the corresponding unitary group. Even when K is noncommutative, we shall again denote by \(\mathrm{T}_{1}(n,\mathrm{K})\) the group of upper triangular matrices of \(\mathbf{M}_{n}(\mathrm{K})\) whose diagonal entries are all 1.

PROPOSITION 7. --- Let \(D_{+}^{*}\) be the group of diagonal matrices with diagonal elements \textgreater0. The mapping \((U,D,T)\mapsto UDT\) is a homeomorphism of \(\mathbf{U}(n,\mathbf{K})\times\mathbf{D}_{+}^{*}\times\mathbf{T}_{1}(n,\mathbf{K})\) onto \(\mathbf{GL}(n,\mathbf{K})\) .

Let \((\varepsilon_1, \ldots, \varepsilon_n)\) be the canonical basis of \(\mathbf{K}^n\) . Let \(X \in \mathbf{GL}(n, \mathbf{K})\) . Then the \(X \cdot \varepsilon_i = f_i\) form a basis of \(\mathbf{E}\) . To this basis \((f_i)\) one can associate a basis \((e_i)\) as in Lemma 10. Let \(U\) be the matrix of the unitary automorphism of \(\mathbf{E}\) that transforms \(\varepsilon_i\) into \(e_i\) . Then

\[
U ^ {- 1} \cdot f _ {i} = \varepsilon_ {1} \alpha_ {i 1} + \varepsilon_ {2} \alpha_ {i 2} + \dots + \varepsilon_ {i} \alpha_ {i i}
\]

with \(\alpha_{ii} > 0\) for \(i = 1,2,\ldots ,n\) . Thus \(X = UC\) , where \(C\) is the matrix

\[
\left( \begin{array}{c c c c} \alpha_ {1 1} & \alpha_ {2 1} & \ldots & \alpha_ {n 1} \\ 0 & \alpha_ {2 2} & \ldots & \alpha_ {n 2} \\ \vdots & \vdots & \ddots & \vdots \\ 0 & 0 & \ldots & \alpha_ {n n} \end{array} \right).
\]

Moreover, U and C depend continuously on X by Lemma 10. On the other hand, formula (8) shows that C may be put in the form DT with \(D \in D_{+}^{*}\) , \(T \in \mathrm{T}_{1}(n, \mathrm{K})\) , D and T depending continuously on C. The uniqueness of the decomposition X = UDT follows from the uniqueness property of Lemma 10.

The homeomorphism of Prop. 7 is called the Iwasawa decomposition of \(\mathbf{GL}(n,\mathbf{K})\) .

The group \(\mathbf{G} = \mathbf{D}_{+}^{*} \cdot \mathbf{T}_{1}(n, \mathbf{K})\) is the set of upper triangular matrices over K whose diagonal elements are \textgreater0. Let us identity the element \((z_{ij})\) of this group with the element

\[
\left(\left(z _ {i i}\right) _ {1 \leqslant i \leqslant n}, \left(z _ {i j}\right) _ {1 \leqslant i <   j \leqslant n}\right) \in \left(\mathbf {R} _ {+} ^ {*}\right) ^ {n} \times \mathrm{K} ^ {n (n - 1) / 2}.
\]

Arguing exactly as in Example 4, one finds as right Haar measure on this group the measure (when \(\mathbf{K} = \mathbf{R}\) )

\[
\left(\prod_ {i = 1} ^ {n} z _ {i i} ^ {- i}\right) \cdot \left(\bigotimes_ {i = 1} ^ {n} d z _ {i i}\right) \otimes \left(\bigotimes_ {i <   j} d z _ {i j}\right).
\]

Then applying Prop. 13 of §2, No.~9, one sees that if \(\mathbf{GL}(n,\mathbf{K})\) is identified with \(\mathbf{U}(n,\mathbf{K})\times\mathbf{G}\) by the mapping \((U,S)\mapsto US\) , a Haar measure on \(\mathbf{GL}(n,\mathbf{K})\) is given by (when \(\mathbf{K}=\mathbf{R}\) )

\[
\left(\prod_ {i = 1} ^ {n} z _ {i i} ^ {- i}\right) \cdot \alpha \otimes \left(\bigotimes_ {i = 1} ^ {n} d z _ {i i}\right) \otimes \left(\bigotimes_ {i <   n} d z _ {i j}\right),\tag{14}
\]

where \(\alpha\) denotes a Haar measure on \(\mathbf{U}(n,\mathbf{K})\) .

Example 8. --- Spaces of hermitian forms.

In this example, K always denotes one of the fields R, C, H. We write \(\delta = \dim_{R} K\) (thus \(\delta = 1, 2\) or 4). A hermitian form \(\Phi\) on the right vector space \(K^{n}\) may be written

\[
\Phi (x, y) = \Phi (x _ {1}, \dots , x _ {n}, y _ {1}, \dots , y _ {n}) = \sum_ {i, j = 1} ^ {n} \overline {{{{x}}}} _ {i} h _ {i j} y _ {j}
\]

with \(h_{ij} = \overline{h}_{ji}\) for all i and j. We denote by H the vector space over R formed by the hermitian matrices of \(\mathbf{M}_{n}(\mathbf{K})\) . The mapping \((h_{ij}) \mapsto \Phi\) is an isomorphism of H onto the vector space of hermitian forms on \(K^{n}\) , by means of which we shall identify these two spaces. Let \(H_{+}^{*} \subset H\) be the set of nondegenerate positive hermitian forms on \(K^{n}\) . The set \(H_{+}^{*}\) is convex in H; for, if \(\Phi_{1}, \Phi_{2}\) are in \(H_{+}^{*}\) and if \(\lambda, \mu\) are two numbers \textgreater0 such that \(\lambda + \mu = 1\) , it is clear that \(\lambda \Phi_{1} + \mu \Phi_{2}\) is a positive hermitian form; on the other hand, if \((\lambda \Phi_{1} + \mu \Phi_{2})(x, x) = 0\) , then \(\Phi_{1}(x, x) = \Phi_{2}(x, x) = 0\) , therefore x = 0, so that \(\lambda \Phi_{1} + \mu \Phi_{2}\) is nondegenerate. Let us now show that \(H_{+}^{*}\) is an open subset of H. Let S be the set of \(x = (x_{1}, \ldots, x_{n}) \in K^{n}\) such that \(x_{1} \overline{x}_{1} + \cdots + x_{n} \overline{x}_{n} = 1\) ; this is a compact subset of \(K^{n}\) ; if \(\Phi \in H_{+}^{*}\) , the function \(x \mapsto \Phi(x, x)\) is continuous and \textgreater0 on S, hence its infimum is \textgreater0; if \(\Phi' \in H\) is sufficiently near \(\Phi\) , it follows that \(\Phi'(x, x) > 0\) for all \(x \in S\) , so that \(\Phi'\) is positive and nondegenerate.

The general linear group \(\mathbf{GL}(n,\mathbf{K})\) operates continuously on the right in \(\mathfrak{H}\) by \((X,\Phi)\mapsto\Phi\circ X\) , that is, by \((X,H)\mapsto{}^{t}\overline{X}\cdot H\cdot X\) , where \(H\) denotes the hermitian matrix corresponding to \(\Phi\) . It is clear that \(\mathfrak{H}_{+}^{*}\) is stable under \(\mathbf{GL}(n,\mathbf{K})\) . More precisely, by Alg., Ch. IX, §6, No.~1, Cor. 1 of Th. 1, \(^{1}\) \(\mathfrak{H}_{+}^{*}\) is the orbit under \(\mathbf{GL}(n,\mathbf{K})\) of the form \(\sum_{i=1}^{n}\overline{x}_{i}y_{i}\) corresponding to the identity matrix \(I_{n}\) . The stabilizer of this form is \(\mathbf{U}(n,\mathbf{K})\) . By Lemma 2 of App. I, \(\mathfrak{H}_{+}^{*}\) may be identified, as a topological homogeneous space, with \(\mathbf{GL}(n,\mathbf{K})/\mathbf{U}(n,\mathbf{K})\) .

For every \(X \in \mathbf{GL}(n, \mathbf{K})\) , let \(\widetilde{X}\) be the automorphism \(H \mapsto {}^t \overline{X} \cdot H \cdot X\) of the real vector space \(\mathfrak{H}\) . If \(\mu\) denotes the Haar measure of the additive group \(\mathfrak{H}\) , one has \(\widetilde{X}^{-1}(\mu) = |\det \widetilde{X}| \cdot \mu\) (§1, No.~10, Cor. 1 of Prop. 15). Let us show that

\[
| \det \widetilde {X} | = | N (X) | ^ {\lambda},\tag{15}
\]

where N denotes the norm in \(\mathbf{M}_{n}(\mathrm{K})\) regarded as an R-algebra, and where \(\lambda = 1 - \frac{\delta - 2}{\delta n}\) . It suffices to verify (15) for X running over a system of generators of \(\mathbf{GL}(n,\mathbf{K})\) , hence (A, II, §10, No.~13, Cor. 2 of Prop. 14) for \(X\) of the following types:

\begin{enumerate}
\def\labelenumi{\alph{enumi})}
\tightlist
\item
  \(X\) is the matrix of a mapping of the form
\end{enumerate}

\[
\left(x _ {1}, \dots , x _ {n}\right) \mapsto \left(x _ {\sigma (1)}, \dots , x _ {\sigma (n)}\right),
\]

where \(\sigma \in \mathfrak{S}_n\) . In this case, a power of \(X\) is equal to 1, therefore \(|\det \widetilde{X}| = |\mathrm{N}(X)| = 1\) .

\begin{enumerate}
\def\labelenumi{\alph{enumi})}
\setcounter{enumi}{1}
\tightlist
\item
  \(X\) is the matrix of a mapping of the form
\end{enumerate}

\[
\left(x _ {1}, \dots , x _ {n}\right) \mapsto \left(a x _ {1}, x _ {2}, \dots , x _ {n}\right).
\]

Then, if \((h_{ij})\in \mathfrak{H}\) , one has \(\widetilde{X} ((h_{ij})) = (h_{ij}')\) with \(h_{11}' = \overline{a} h_{11}a = |a|^2 h_{11}\) , \(h_{1i}' = \overline{a} h_{1i}\) for \(i > 1\) , \(h_{ij}' = h_{ij}\) for \(i > 1\) , \(j > 1\) ; therefore

\[
| \det \widetilde {X} | = | a | ^ {2} | a | ^ {\delta (n - 1)} = | a | ^ {2 + \delta (n - 1)}.
\]

On the other hand, if \(Y = (y_{ij}) \in \mathbf{M}_n(\mathbf{K})\) , then \(XY = (y_{ij}')\) with \(y_{1j}' = ay_{1j}\) and \(y_{ij}' = y_{ij}\) for \(i > 1\) ; therefore \(|\mathrm{N}(X)| = |a|^{\delta n}\) . The formula (15) is again verified.

\begin{enumerate}
\def\labelenumi{\alph{enumi})}
\setcounter{enumi}{2}
\tightlist
\item
  \(X\) is the matrix of a mapping of the form
\end{enumerate}

\[
\left(x _ {1}, \dots , x _ {n}\right) \mapsto \left(x _ {1} + b x _ {2}, x _ {2}, \dots , x _ {n}\right).
\]

Then \(\widetilde{X}\big((h_{ij})\big) = (h_{ij}^{\prime})\) with \(h_{11}^{\prime} = h_{11}\) , \(h_{12}^{\prime} = h_{12} + h_{11}b\) , \(h_{1i}^{\prime} = h_{1i}\) for \(i > 2\) , \(h_{22}^{\prime} = h_{22} + \overline{b}h_{12} + \overline{h}_{12}b + \overline{b}h_{11}b\) , \(h_{2i}^{\prime} = h_{2i} + \overline{b}h_{1i}\) for \(i > 2\) , \(h_{ij}^{\prime} = h_{ij}\) for \(i > 2\) , \(j > 2\) . Taking into account Lemma 6, one sees that \(|\det \widetilde{X}| = 1\) . One verifies similarly that \(|\mathrm{N}(X)| = 1\) , which completes the proof of formula (15).

This established, the measure \(|\mathrm{N}(H)|^{-\lambda/2}d\mu(H)\) on H is invariant under \(\mathbf{GL}(n,\mathbf{K})\) , since

\[
\begin{array}{r l} & {\widetilde {X} ^ {- 1} \big (| \mathrm{N} (H) | ^ {- \lambda / 2} d \mu (H) \big) = \big | \mathrm{N} \big (\widetilde {X} (H) \big) \big | ^ {- \lambda / 2} \cdot | \det \widetilde {X} | d \mu (H)} \\ & {\qquad = | \mathrm{N} (H) | ^ {- \lambda / 2} | \mathrm{N} (X) | ^ {- \lambda} | \det \widetilde {X} | d \mu (H) = | \mathrm{N} (H) | ^ {- \lambda / 2} d \mu (H).} \end{array}
\]

If \(H \in \mathfrak{H}_+^*\) , then \(H = \widetilde{X}(\mathrm{I}_n) = {}^t\overline{X}X\) for some \(X \in \mathbf{GL}(n,\mathbf{K})\) , therefore \(\mathrm{N}(H) = \overline{\mathrm{N}(X)}\mathrm{N}(X) > 0\) . Consequently, on \(\mathfrak{H}_+^*\) , the unique (up to a constant factor) measure invariant under \(\mathbf{GL}(n,\mathbf{K})\) (cf.~§2, No.~6, Th. 3) is the measure

\[
d \gamma (H) = \mathrm{N} (H) ^ {- \lambda / 2} d \mu (H).
\]

In particular,

\[
d \gamma (H) = (\det H) ^ {- (n + 1) / 2} d \mu (H) \quad \text { when } K = R,
\]

\[
d \gamma (H) = (\det H) ^ {- n} d \mu (H) \quad \text { when } \mathrm{K} = \mathrm{C}.
\]

\appendixheading

*Lemma 1. --- Let X be a locally compact space, R an open equivalence relation in X, such that the quotient space X/R is paracompact; let \(\pi\) be the canonical mapping of X onto X/R. There exists a continuous real-valued function \(F \geqslant 0\) on X such that:

\begin{enumerate}
\def\labelenumi{\alph{enumi})}
\item
  F is not identically zero on any equivalence class with respect to R;
\item
  for every compact subset K of X/R, the intersection of \(\overline{\pi}^{1}(K)\) with Supp F is compact.
\end{enumerate}

To each point \(u \in \mathbf{X} / \mathbf{R}\) let us associate a function \(f_{u} \in \mathcal{K}_{+}(\mathbf{X})\) such that \(f_{u}\) is not identically zero on \(\overline{\pi}(u)\) ; let \(\Omega_{u}\) be the open set of points where \(f_{u} > 0\) ; thus \(u \in \pi(\Omega_{u})\) . Since \(\pi\) is an open mapping, the \(\pi(\Omega_{u})\) form an open covering of \(\mathbf{X} / \mathbf{R}\) . There exists a locally finite open covering \((\mathrm{U}_{\iota})_{\iota \in \mathbf{I}}\) , finer than the covering by the \(\pi(\Omega_{u})\) , then (GT, IX, §4, No.~3, Prop. 3) a partition of unity \((g_{\iota})_{\iota \in \mathbf{I}}\) on \(\mathbf{X} / \mathbf{R}\) subordinate to the covering \((\mathrm{U}_{\iota})\) . For every \(\iota \in \mathbf{I}\) , choose a \(u_{\iota}\) such that \(\mathrm{U}_{\iota} \subset \pi(\Omega_{u_{\iota}})\) . The function \(\mathrm{F}_{\iota} = (g_{\iota} \circ \pi) \cdot f_{u_{\iota}}\) belongs to \(\mathcal{K}(\mathbf{X})\) and has support contained in \(\overline{\pi}(\mathrm{U}_{\iota})\) . The supports of the \(F_{\iota}\) therefore form a locally finite family, so that one can define a continuous function \(F \geqslant 0\) on \(X\) by setting \(F = \sum_{\iota \in I} F_{\iota}\) .

For every \(u \in X/R\) , there exists an \(\iota\) such that \(g_{\iota}(u) > 0\) and therefore \(u \in U_{\iota}\) ; next, there exists an \(x \in \Omega_{u_{\iota}}\) such that \(\pi(x) = u\) ; then \(f_{u_{\iota}}(x) > 0\) and \(g_{\iota}(\pi(x)) > 0\) , therefore \(F_{\iota}(x) > 0\) and a fortiori \(F(x) > 0\) ; this proves that F has the property a). Finally, let K be a compact subset of X/R. There exists a finite subset J of I such that, for \(\iota \in I - J\) , one has \(\mathrm{U}_{\iota}\cap \mathrm{K} = \varnothing\) , therefore \(\overline{\pi} (\mathbf{K})\cap \operatorname {Supp}\mathbf{F}_{\iota} = \varnothing\) . Then the set

\[
\overline {{\pi}} ^ {1} (\mathrm{K}) \cap \operatorname{Supp} \mathrm{F} = \overline {{\pi}} ^ {1} (\mathrm{K}) \cap \left(\bigcup_ {\iota \in \mathrm{I}} \operatorname{Supp} \mathrm{F} _ {\iota}\right) = \overline {{\pi}} ^ {1} (\mathrm{K}) \cap \left(\bigcup_ {\iota \in \mathrm{J}} \operatorname{Supp} \mathrm{F} _ {\iota}\right)
\]

is compact.

Lemma 2. --- Let G be a locally compact group countable at infinity, M a Baire space. Suppose that G operates on the left continuously and transitively in M. For every \(x \in M\) , let \(H_{x}\) be the stabilizer of x in G, so that the mapping \(s \mapsto sx\) of G onto M defines, by passage to the quotient, a continuous bijection \(\varphi_{x}\) of \(G/H_{x}\) onto M. Then \(\varphi_{x}\) is a homeomorphism of \(G/H_{x}\) onto M (in other words (GT, III, §2, No.~5) M is a topological homogeneous space).

Let \(x_0 \in \mathbf{M}\) . It suffices to prove (loc. cit., Prop. 15) that the mapping \(s \mapsto sx_0\) transforms every neighborhood \(\mathbf{V}\) of \(e\) in \(\mathbf{G}\) into a neighborhood of \(x_0\) in \(\mathbf{M}\) . Let \(\mathbf{W}\) be a symmetric compact neighborhood of \(e\) such that \(\mathbf{W}^2 \subset \mathbf{V}\) . By hypothesis, \(\mathbf{G}\) is the union of a sequence of compact sets, hence of a sequence \((s_n\mathbf{W})\) of translates of \(\mathbf{W}\) . Then \(\mathbf{M}\) is the union of the sequence of compact sets \((s_n\mathbf{W}x_0)\) . Since \(\mathbf{M}\) is a Baire space, there exists an index \(n\) such that \(s_n\mathbf{W}x_0\) has an interior point \(s_nwx_0\) (where \(w \in \mathbf{W}\) ). Consequently \(x_0\) is an interior point of

\[
w ^ {- 1} s _ {n} ^ {- 1} (s _ {n} \mathrm{W} x _ {0}) = w ^ {- 1} \mathrm{W} x _ {0} \subset \mathrm{V} x _ {0},
\]

so that \(Vx_{0}\) is a neighborhood of \(x_{0}\) in M.

\appendixheading

*Lemma 1. --- Let X, B be two locally compact spaces, \(\pi\) a mapping of X into B, \(\nu\) a positive measure on B. Let \(b \mapsto \lambda_b\) ( \(b \in \mathbf{B}\) ) be a \(\nu\) -adequate family of positive measures on X such that, for every \(b \in \mathbf{B}\) , the measure \(\lambda_b\) is concentrated on \(\overline{\pi}^1(b)\) . Set \(\mu = \int \lambda_b d\nu(b)\) , and assume that the mapping \(\pi\) is \(\mu\) -measurable.

\begin{enumerate}
\def\labelenumi{\alph{enumi})}
\item
  If \(\mathbf{N} \subset \mathbf{B}\) is locally \(\nu\) -negligible, then \(\overline{\pi}^1(\mathbf{N})\) is locally \(\mu\) -negligible.
\item
  If \(f\) is a \(\nu\) -measurable function on \(\mathbf{B}\) (with values in a topological space), then \(f \circ \pi\) is \(\mu\) -measurable on \(\mathbf{X}\) .
\end{enumerate}

Let K be a compact subset of X. We are to prove that \(\overline{\pi}^{1}(N) \cap K\) is \(\mu\) -negligible and that the restriction of \(f \circ \pi\) to K is \(\mu\) -measurable. Now, K is the union of a \(\mu\) -negligible set and a sequence of compact subsets \(K_{n}\) such that \(\pi|K_{n}\) is continuous. It suffices to show that \(\overline{\pi}^{1}(N) \cap K_{n}\) is \(\mu\) -negligible and that the restriction of \(f \circ \pi\) to \(K_{n}\) is \(\mu\) -measurable. We may therefore assume henceforth that \(\pi|K\) is continuous. Then \(\pi(K) = K'\) is compact. Since \(\overline{\pi}^{1}(N) \cap K = \overline{\pi}^{1}(N \cap K') \cap K\) and \(N \cap K'\) is \(\nu\) -negligible, we may assume henceforth that N is \(\nu\) -negligible. Then N is contained in a \(\nu\) -negligible set \(N'\) that is a countable intersection of open sets (Ch. IV, §4, No.~6, Cor. 2 of Th. 4). Since \(\pi\) is \(\mu\) -measurable, \(\overline{\pi}^{1}(N')\) is \(\mu\) -measurable (Ch. IV, §5, No.~5, Prop. 7), therefore \(\overline{\pi}^{1}(N') \cap K\) is \(\mu\) -integrable, and (Ch. V, §3, No.~4, Th. 1)

\[
\mu \big (\overline {{{{\pi}}}} ^ {1} (\mathbf {N} ^ {\prime}) \cap \mathbf {K} \big) = \int_ {\mathbf {B}} \lambda_ {b} \big (\overline {{{{\pi}}}} ^ {1} (\mathbf {N} ^ {\prime}) \cap \mathbf {K} \big) d \nu (b).
\]

Now, if \(b \notin \mathbf{N}'\) , \(\overline{\pi}(\mathbf{N}') \cap \mathbf{K}\) is \(\lambda_b\) -negligible since \(\lambda_b\) is concentrated on \(\overline{\pi}(b)\) by hypothesis. Therefore \(\mu\big(\overline{\pi}(\mathbf{N}') \cap \mathbf{K}\big) = 0\) . A fortiori, \(\overline{\pi}(\mathbf{N}) \cap \mathbf{K}\) is \(\mu\) -negligible and a) has indeed been proved. On the other hand, there exists a partition of \(\mathbf{K}'\) formed by a \(\nu\) -negligible set \(\mathbf{M}\) and a sequence \((\mathbf{K}_n')\) of compact sets such that \(f|_{\mathbf{K}_n'}\) is continuous. Then the restriction of \(f \circ \pi\) to each set \(\overline{\pi}^1(\mathbf{K}) \cap \mathbf{K}_n'\) is continuous, and \(\overline{\pi}^1(\mathbf{M}) \cap \mathbf{K}\) is \(\mu\) -negligible by a), therefore the restriction of \(f \circ \pi\) to \(\mathbf{K}\) is indeed \(\mu\) -measurable.

\specialchapter{Exercises}

\exercisesection{}

\begin{enumerate}
\def\labelenumi{\arabic{enumi})}
\item
  Let \(G\) be a compact group. Show that every continuous representation \(\varphi\) of \(G\) in \(\mathbf{R}_+^*\) is such that \(\varphi(G) = \{1\}\) . From this, deduce that a relatively invariant positive measure on \(G\) is invariant.
\item
  Let G be a topological group, such that the commutator subgroup of G is dense in G. Show that every continuous representation \(\varphi\) of G in a Hausdorff abelian group satisfies \(\varphi(G) = \{e\}\) . From this, deduce that if G is locally compact, then every relatively invariant complex measure on G is invariant. In particular, G is unimodular.
\item
  Let \(G\) be a locally compact group, \(\mu\) a left Haar measure on \(G\) , and \(\nu = \Delta_G^{-1/2} \cdot \mu\) . Show that
\end{enumerate}

\[
\boldsymbol {\gamma} (s) \nu = \Delta_ {\mathrm{G}} (s) ^ {1 / 2} \nu , \quad \boldsymbol {\delta} (s) \nu = \Delta_ {\mathrm{G}} (s) ^ {1 / 2} \nu , \quad \stackrel {\vee} {\nu} = \nu .
\]

\begin{enumerate}
\def\labelenumi{\arabic{enumi})}
\setcounter{enumi}{3}
\item
  Let \(G, G'\) be two locally compact groups, \(V\) (resp. \(V'\) ) an open neighborhood of the neutral element of \(G\) (resp. \(G'\) ), \(\varphi\) a local isomorphism of \(G'\) with \(G\) , defined on \(V'\) , such that \(\varphi(V') = V\) . Show that \(\Delta_G \circ \varphi\) is the restriction of \(\Delta_{G'}\) to \(V'\) .
\item
  For every \(a\) belonging to the multiplicative group \(\mathbf{Q}^*\) , let \(\varphi(a)\) be the automorphism of the additive group \(\mathbf{R}\) defined by \(\varphi(a)x = ax\) . Equip \(\mathbf{Q}^*\) with the discrete topology. Let \(\mathbf{G}\) be the topological semi-direct product of \(\mathbf{Q}^*\) and \(\mathbf{R}\) defined by \(\varphi\) (GT, III, §2, No.~10). Show that \(\mathbf{G}\) is locally isomorphic to \(\mathbf{R}\) , but is not unimodular.
\item
  Let \(\mathbf{G}\) be a locally compact group having an open and compact subgroup \(\mathbf{H}\) . Show that for every automorphism \(\varphi\) of \(\mathbf{G}\) , mod \(\varphi \in \mathbf{Q}_+^*\) . (Observe that \(\varphi(\mathbf{H}) \cap \mathbf{H}\) has finite index in \(\mathbf{H}\) and in \(\varphi(\mathbf{H})\) .) Show that the equality \(\Delta_{\mathbf{G}}(s) = 1\) defines an open subgroup of \(\mathbf{G}\) containing \(\mathbf{H}\) .
\item
  Let G be a locally compact group, \(\beta\) a left Haar measure on G, A a subset of G, and B a relatively compact \(\beta\) -integrable subset of G such that \(\beta(\mathrm{B}) > 0\) . Show that if the compact subsets of \(\mathbf{A} \cdot \mathbf{B}\) have bounded measures for \(\beta\) , then \(\mathbf{A}\) is relatively compact. (Imitate the argument of Prop. 2.)
\item
  Let G be a locally compact group, A a dense subset of G, \(\alpha\) a left Haar measure on G, and H an \(\alpha\) -measurable subset of G having the following property: for every \(s \in A\) , \(sH \cap (\mathbb{C}H)\) and \(H \cap (\mathbb{C}sH)\) are locally \(\alpha\) -negligible. Show that either H is locally \(\alpha\) -negligible or its complement is locally \(\alpha\) -negligible. (Show that \(\varphi_{H} \cdot \alpha\) is left invariant.)
\item
  Adopt the notations of the proof of Th. 1 of No.~2. Let \(\beta\) be the unique left Haar measure on G such that \(\beta(f_0) = 1\) . Show that, for every \(f \in \mathcal{K}(\mathbf{G})\) , \(\mathrm{I}_g(f)\) tends to \(\beta(f)\) with respect to \(\mathfrak{B}\) . (Let \(a\) be a cluster point of \(g \mapsto \mathrm{I}_g(f)\) with respect to \(\mathfrak{B}\) . There exists an ultrafilter \(\mathfrak{U}\) finer than \(\mathfrak{B}\) such that \(\mathrm{I}_g(f)\) tends to \(a\) with respect to \(\mathfrak{U}\) . On the other hand, \(\mathrm{I}_g(f)\) tends to \(\beta(f)\) with respect to \(\mathfrak{U}\) .)
\item
  Let G be a locally compact group, \(\mu\) a left Haar measure on G. Show that every \(\mu\) -integrable set A is contained in a countable union of compact sets. (Reduce to the case that A is open, and observe that there exists an open subgroup of G that is a countable union of compact subsets.)
\end{enumerate}

¶ 11) Let G be a nondiscrete locally compact group countable at infinity, β a left Haar measure on G. Show that every compact neighborhood V of e contains a normal subgroup H of G, β-negligible and such that G/H is a Polish locally compact group. (One can proceed as follows: let L be the open subgroup of G generated by V; let \(b_{1}, b_{2}, \ldots\) be representatives of the left cosets with respect to L. Form a decreasing sequence \((\mathrm{V}_{n})\) of symmetric neighborhoods of e such that: 1) \(\mathrm{V}_{n}^{2} \subset \mathrm{V}_{n-1}\) ; 2) \(x\mathrm{V}_{n}x^{-1} \subset \mathrm{V}_{n-1}\) for every \(x \in \mathrm{V}\) ; 3) \(b_{i}\mathrm{V}_{n}b_{i}^{-1} \subset \mathrm{V}_{n-1}\) for \(1 \leqslant i \leqslant n\) ; 4) \(\beta(\mathrm{V}_{n}) \leqslant 1/n\) . Set \(\mathrm{H} = \mathrm{V}_{1} \cap \mathrm{V}_{2} \cap \cdots\) .

Let \((f_{n})\) be a sequence of numerical functions uniformly continuous for the left uniform structure of G. Show that H can be constructed in such a way that, in addition, the \(f_{n}\) are constant on the cosets of H. (In the preceding construction, take \(V_{n}\) to be such that \(|f_{i}(x) - f_{i}(y)| \leqslant 1/n\) for \(x^{-1}y \in V_{n}\) and \(1 \leqslant i \leqslant n\) .)

Let \((g_n)\) be a sequence of \(\beta\) -integrable numerical functions on \(G\) . Show that \(H\) can be constructed in such a way that, in addition, each \(g_n\) is equal almost everywhere to a function that is constant on the cosets of \(H\) .

\begin{enumerate}
\def\labelenumi{\arabic{enumi})}
\setcounter{enumi}{11}
\tightlist
\item
  Let \(\mathbf{K}\) be a nondiscrete locally compact field, \(\mathbf{E}\) a left topological vector space over \(\mathbf{K}\) .
\end{enumerate}

\begin{enumerate}
\def\labelenumi{\alph{enumi})}
\item
  If E is 1-dimensional, then E is isomorphic to \(\mathbf{K}_s\) (imitate the proof of Prop. 2 of TVS, I, §2, No.~2, replacing the absolute value by the function mod \(_K\) ).
\item
  If E is n-dimensional, then E is isomorphic to \(K_{s}^{n}\) (imitate the proof of Th. 2 of TVS, I, §2, No.~3).
\item
  Assume that E is locally compact. Let F be a linear subspace of E of finite dimension n.~For every \(a \in K\) , let \(\text{mod}_{\mathbf{E}}(a)\) be the modulus of \(x \mapsto ax\) in E. Since F is closed in E by b), one can form \(\text{mod}_{\mathbf{E}/\mathbf{F}}(a)\) . Show that \(\text{mod}_{\mathbf{E}}(a) = \text{mod}_{\mathbf{K}}(a)^n \text{mod}_{\mathbf{E}/\mathbf{F}}(a)\) . Show that if \(\text{mod}_{\mathbf{K}}(a) < 1\) and \(F \neq E\) , then \(\text{mod}_{\mathbf{E}/\mathbf{F}}(a) < 1\) . (One has \(a^n \to 0\) as \(n \to +\infty\) ; from this, deduce that \(\text{mod}_{\mathbf{E}/\mathbf{F}}(a^n) \to 0\) by arguing as in Prop. 12 of No.~10.) Deduce from this that \(n \leqslant \log \text{mod}_{\mathbf{E}}(1/a)/\log \text{mod}_{\mathbf{K}}(1/a)\) , hence that E is finite-dimensional.
\end{enumerate}

(We shall see in CA, VI, that the topology of a locally compact field can be defined by an absolute value. The final result of c) will then be a special case of Th. 3 of TVS, I, §2, No.~4.)

¶ 13) Let G be a locally compact group operating continuously on the left in a locally compact Polish space T, β a left Haar measure on G, and ν a positive measure on T quasi-invariant under G.

\begin{enumerate}
\def\labelenumi{\alph{enumi})}
\tightlist
\item
  There exists a function \((s,x)\mapsto \chi (s,x) > 0\) on \(\mathbf{G}\times \mathbf{T}\) , locally \((\beta \otimes \nu)\) -integrable, such that for every \(s\in \mathbf{G}\) , the function \(x\mapsto \chi (s^{-1},x)\) is locally \(\nu\) -integrable and satisfies \(\gamma(s)\nu = \chi(s^{-1},\cdot)\cdot\nu\) . (Show that \((s,x)\mapsto(s,sx)\) transforms \(\beta\otimes\nu\) into an equivalent measure, by making use of criterion 3) of Th. 2 of Ch. V, §5, No.~5; let \(\chi(s^{-1},x)d\beta(s)d\nu(x)\) be this equivalent measure. For \(\varphi\in\mathcal{H}(\mathrm{G})\) and \(\psi\in\mathcal{H}(\mathrm{T})\) ,
\end{enumerate}

\[
\int \varphi (s) d \beta (s) \int \psi (s x) d \nu (x) = \int \varphi (s) d \beta (s) \int \psi (x) \chi (s ^ {- 1}, x) d \nu (x),
\]

whence \(\int \psi(sx) d\nu(x) = \int \psi(x)\chi(s^{-1}, x) d\nu(x)\) except on a locally \(\beta\) -negligible set \(N(\psi)\) . Next make use of Lemma 1 of Ch. VI, §3, No.~1; then modify \(\chi\) on a \((\beta \otimes \nu)\) -negligible set.)

\begin{enumerate}
\def\labelenumi{\alph{enumi})}
\setcounter{enumi}{1}
\tightlist
\item
  Show that for any \(s, t\) in \(\mathbf{G}\) , one has
\end{enumerate}

\[
\chi (s t, x) = \chi (s, t x) \chi (t, x)
\]

except on a \(\nu\) -negligible set of values of \(x\) (make use of the relation \(\pmb{\gamma}(st)\nu = \pmb{\gamma}(s)(\pmb{\gamma}(t)\nu)\) ). c) Show that the function \(\chi\) of \(a\) ) is determined up to a locally \((\beta \otimes \nu)\) -negligible subset of \(\mathbf{G} \times \mathbf{T}\) .

¶ 14) Let T be a locally compact space, U a uniform structure on T for which there exists an entourage \(V_{0}\) such that \(V_{0}(t)\) is compact for all \(t \in T\) (GT, II, §4, Exer. 9). On the other hand, let \(\Gamma\) be a uniformly equicontinuous group of homeomorphisms of T (for the uniform structure U), such that there exists an \(a \in T\) whose orbit under \(\Gamma\) is dense. Show that there exists on T a positive measure \(\neq 0\) invariant under \(\Gamma\) , and that this measure is unique up to a constant factor. One may proceed as follows:

\(1^{\circ}\) As for the existence of the invariant measure, follow the method of No.~2, Th. 1; one proves first that if K is a compact subset of T and U is an open neighborhood of \(a\) , there exists a finite number of elements \(\sigma_{i} \in \Gamma\) such that \(\mathrm{K} \subset \bigcup \sigma_{i}(\mathrm{U})\) .

\(2^{\circ}\) As for uniqueness, observe first that the entourages of \(\mathcal{U}\) that are invariant under every homeomorphism \(\sigma \times \sigma\) of T × T, for \(\sigma \in \Gamma\) , form a fundamental system of entourages \(\mathfrak{S}\) . For every relatively compact set A ⊂ T and every relatively compact set B ⊂ T with nonempty interior, let (A : B) be the smallest number of elements of a covering of A by sets of the form \(\sigma\) B, where \(\sigma \in \Gamma\) ; if C ⊂ T is a third relatively compact set with nonempty interior, then (A : C) ≤ (A : B)(B : C). Let K be a compact subset of T, L ⊃ K a relatively compact open set, V ∈ \(\mathfrak{S}\) a symmetric open entourage contained in V0 such that V(K) ⊂ L, and W ∈ \(\mathfrak{S}\) a symmetric closed entourage contained in V. For every symmetric entourage U ∈ \(\mathfrak{S}\) such that UW ⊂ V and WU ⊂ V, show that, for every positive measure \(\nu\) on T invariant under \(\Gamma\) ,

\begin{enumerate}
\def\labelenumi{(\arabic{enumi})}
\tightlist
\item
\end{enumerate}

\[
\big (\mathrm{W} (a): \mathrm{U} (a) \big) \nu (\mathrm{K}) \leqslant \big (\mathrm{L}: \mathrm{U} (a) \big) \nu \big (\mathrm{V} (a) \big)\tag{2}
\]

\[
\bigl (\mathrm{K}: \mathrm{U} (a) \bigr) \nu \bigl (\mathrm{W} (a) \bigr) \leqslant \bigl (\mathrm{V} (a): \mathrm{U} (a) \bigr) \nu (\mathrm{L}).
\]

For this, one proves that if \((\sigma_i(\mathrm{U}(a)))\) is a covering of L formed by \((\mathrm{L}:\mathrm{U}(a))\) sets, every \(x\in \mathbf{K}\) belongs to at least \((\mathrm{W}(a):\mathrm{U}(a))\) sets of the form \(\sigma_{i}(\mathrm{V}(a))\) , and that every \(y\in \mathrm{L}\) belongs to at most \((\mathrm{V}(a):\mathrm{U}(a))\) sets of the form \(\sigma_{i}(\mathrm{W}(a))\) (note that if \(y\in \sigma_i(\mathrm{W}(a))\) , then \(\sigma_{i}(\mathrm{U}(a))\) is contained in \(\mathrm{V}(x)\) ).

Let \(\mathfrak{F}\) be an ultrafilter finer than the section filter of \(\mathfrak{S}\) ; let \(A_0\) be a nonempty, relatively compact open set in T, and set \(\lambda_U(A) = (A: U(a)) / (A_0: U(a))\) for every symmetric entourage \(U \in \mathfrak{S}\) and every relatively compact subset A of T. Let \(\lambda(A) = \lim_{\mathfrak{F}} \lambda_U(A)\) ; for every compact subset K of T, set \(\lambda'(K) = \inf \lambda(B)\) , where B runs over the set of relatively compact open neighborhoods of K. Deduce from (1) and (2) that

\[
\begin{array}{r} \lambda^ {\prime} \big (\mathrm{W} (a) \big) \nu (\mathrm{K}) \leqslant \lambda (\mathrm{L}) \nu \big (\mathrm{W} (a) \big) \\ \lambda^ {\prime} (\mathrm{K}) \nu \big (\mathrm{W} (a) \big) \leqslant \lambda^ {\prime} \big (\mathrm{W} (a) \big) \nu (\mathrm{L}). \end{array}
\]

Conclude from this that if \(K_{1}, K_{2}\) are two compact subsets of T, and \(L_{1} \supset K_{1}\) , \(L_{2} \supset K_{2}\) are two relatively compact open sets, then

\[
\lambda^ {\prime} (\mathrm{K} _ {1}) \nu (\mathrm{K} _ {2}) \leqslant \lambda (\mathrm{L} _ {2}) \nu (\mathrm{L} _ {1}),
\]

and finally that \(\lambda^{\prime}(\mathrm{K}_{1})\nu (\mathrm{K}_{2}) = \lambda^{\prime}(\mathrm{K}_{2})\nu (\mathrm{K}_{1})\)

¶ 15) Let G be a locally compact group, H a closed subgroup of G, so that G operates continuously on G/H. Assume that H satisfies the following condition: for every neighborhood V of e in G, there exists a neighborhood U of e such that HU ⊂ VH (note that this condition is satisfied for every subgroup H of G if e admits a fundamental system of neighborhoods invariant under every inner automorphism of G). Show that there then exists a nonzero positive measure on G/H invariant under G. (Use the same method as in the proof of Th. 1 of No.~2, on observing, on the one hand, that as U runs over the set of neighborhoods of e in G, the images of the sets HUH under the canonical mapping π : G → G/H form a fundamental system of neighborhoods of π(H); on the other hand, for every neighborhood V of e in G and every compact subset K of G, there exist a neighborhood U of e in G and a compact subset L of G such that every set of the form sHUH that intersects KH is of the form s'HUH with s' ∈ L.)

¶ 16) Let G be a compact abelian group, E a dense subset of G that is stable under the law of composition of G. Assume that, for every bounded numerical function f on E, there is defined a number M(f) having the following property: for any two bounded functions f, g on E, for any elements \(a_{i}, b_{k}\) of E ( \(1 \leqslant i \leqslant m, 1 \leqslant k \leqslant n\) ), and for any real numbers \(\alpha_{i}, \beta_{k}\) ( \(1 \leqslant i \leqslant m, 1 \leqslant k \leqslant n\) ) such that

\[
\sum_ {i} \alpha_ {i} f (x a _ {i}) \leqslant \sum_ {k} \beta_ {k} g (x b _ {k})
\]

for all \(x \in \mathbf{E}\) , one then has the inequality

\[
\left(\sum_ {i} \alpha_ {i}\right) \mathrm{M} (f) \leqslant \left(\sum_ {k} \beta_ {k}\right) \mathrm{M} (g).
\]

Assume in addition that \(\mathbf{M}(1) = 1\) (cf.~TVS, IV, Appendix).

\begin{enumerate}
\def\labelenumi{\alph{enumi})}
\item
  Let \(\mu\) be the normalized Haar measure on G. Show that for every continuous numerical function \(f\) on G, \(\int f d\mu = \mathrm{M}(f|\mathrm{E})\) . (By considering a suitable partition of unity, approximate the constant function equal to \(\int f d\mu\) on E by means of functions of the form \(x \mapsto \sum_{i} \alpha_i f(xa_i)\) , where \(a_i \in \mathrm{E}\) .)
\item
  Deduce from a) that if one sets \(\nu(B) = M(\varphi_B)\) for every subset B of E, then \(\nu(U \cap E) \geqslant \mu(U)\) for every open subset U of G, and \(\nu(F \cap E) \leqslant \mu(F)\) for every closed subset F of G. For every \(\mu\) -quadratic subset P of G (Ch. IV, §5, Exer. 17 d)), one has \(\nu(P \cap E) = \mu(P)\) .
\end{enumerate}

¶ 17) Let T be a locally compact space, S a subset of T equipped with a law of composition \((x,y)\mapsto xy\) that makes it a monoid \(^{1}\) (not necessarily having a priori a neutral element), and which is continuous on S×S when S is equipped with the topology induced by that of T. Assume in addition that every element of S is cancellable \(^{2}\) (A, I, §2, No.~2). Let \(\mu\) be a bounded positive measure on T, concentrated on S (so that S is \(\mu\) -measurable) and of total mass 1. Assume that \(\mu\) is left invariant in the following sense: for every numerical function f defined on S, continuous and bounded (hence \(\mu\) -measurable by Ch. IV, §5, No.~5, Cor. of Prop. 8), one has \(\int f(sx) d\mu(x) = \int f(x)d\mu(x)\) for all \(s \in S\) . It comes to the same to say that \(\mu(sK) = \mu(K)\) for every compact subset K of S.

\begin{enumerate}
\def\labelenumi{\alph{enumi})}
\item
  Consider the product measure \(\mu\otimes\mu\) on T×T; show that there exist compact subsets K of S such that the images of K×K under the two continuous mappings \((x,y)\mapsto(x,xy)\) and \((x,y)\mapsto(xy,x)\) have measure arbitrarily close to 1 (use the Lebesgue--Fubini theorem). Conclude from this that these images have a common point, and deduce therefrom that S has a neutral element (cf.~A, I, §2, Exer. 9).
\item
  Show that S is a compact group. (First prove that for every \(x \in S\) , the inner measure \(\mu_{*}(xS)\) is equal to 1, hence that xS is measurable and that the measure is concentrated on xS; xS is then a monoid to which the result of a) can be applied, which shows that S is a group. Argue as in Prop. 2 to prove that S is compact. Finally, make use of Exer. 21 of GT, III, §4.3)
\end{enumerate}

18) Let X be a locally compact space, G a compact group operating continuously on X, E the orbit of a point of X, and \(\mathcal {F}\) a vector space of continuous numerical functions on X such that, for every function \(f \in \mathcal {F}\) and every \(s \in G\) , one has \(\gamma (s) f \in \mathcal {F}\) ; assume in addition that \(\mathcal {F}\) contains the constant functions on X. Let \(x _ {0} \in \mathrm{X}\) be a point invariant under G and such that \(| f (x _ {0}) | \leqslant \sup _ {y \in E} | f (y) |\) for every function \(f \in \mathcal {F}\) . Show

that one then has \(f(x_0) = \int_{\mathbf{G}} f(s \cdot z) ds\) for every \(z \in \mathbf{E}\) and every \(f \in \mathcal{F}\) . (Show that there exists a positive measure \(\nu\) on \(\mathbf{E}\) , of total mass 1, such that \(f(x_0) = \int_{\mathbf{E}} f(z) d\nu(z)\) for every \(f \in \mathcal{F}\) , and apply the Lebesgue-Fubini theorem.) *The case that \(X = R^n\) , \(G\) is the orthogonal group, and \(x_0 = 0\) ; apply the formula (7) of §2.*

\begin{enumerate}
\def\labelenumi{\arabic{enumi})}
\setcounter{enumi}{18}
\tightlist
\item
  Let X be a compact space, A a normed algebra over R with unity element, and G a compact group; assume that G operates continuously on A and on X and that, for every \(s \in G\) , \(a \mapsto s \cdot a\) is an automorphism of the algebra A such that \(\|s \cdot a\| = \|a\|\) for all \(a \in A\) . A mapping f of X into A is said to be covariant under G if
\end{enumerate}

\[
f (s \cdot x) = s \cdot f (x)
\]

for all \(s \in G\) and \(x \in X\) . Let \(B\) be a subring of \(A^X\) consisting of covariant continuous functions and containing all of the covariant continuous mappings of \(X\) into \(R\) ( \(R\) being identified with a subalgebra of \(A\) ; one observes that for such a mapping \(g\) , \(g(s \cdot x) = g(x)\) for all \(s \in G\) and \(x \in X\) ). Let \(f\) be a covariant continuous mapping of \(X\) into \(A\) , and assume that for every \(y \in X\) , there exists a mapping \(g_y \in B\) such that \(f(y) = g_y(y)\) . Then, for every \(\varepsilon > 0\) , there exists a mapping \(g \in B\) such that \(\|f(x) - g(x)\| \leqslant \varepsilon\) for all \(x \in X\) . (Make use of a suitable continuous partition of unity ( \(\varphi_i\) ) on \(X\) and introduce the functions \(h_i\) such that \(h_i(x) = \int_G \varphi_i(s \cdot x) ds\) .)

¶ 20) Let G be a locally compact group, μ a left Haar measure on G, A and B two subsets of G.

\begin{enumerate}
\def\labelenumi{\alph{enumi})}
\tightlist
\item
  Assume that one of the following two conditions holds:
\end{enumerate}

\(\alpha)\) A is \(\mu\) -integrable;

\(\beta)\ \mu^{*}(A) < +\infty\) and \(B\) is \(\mu\) -measurable.

Show that, in each of these two cases, the function \(f(s) = \mu^{*}(sA \cap B)\) is uniformly continuous on \(G\) for the right uniform structure of \(G\) . (For any two subsets \(M, N\) of \(G\) , one sets

\[
d (\mathrm{M}, \mathrm{N}) = \mu^ {*} \big ((\mathrm{M} \cap \mathbb {C N}) \cup (\mathrm{N} \cap \mathbb {C M}) \big).
\]

First consider the case that A is compact; making use of Ch. IV, §4, No.~6, Th. 4, show that for every \(\varepsilon > 0\) , there exists a neighborhood U of \(e\) in G such that, for every \(s \in G\) and \(t \in U\) ,

\[
d (s \mathrm{A} \cap \mathrm{B}, s t \mathrm{A} \cap \mathrm{B}) \leqslant \varepsilon .
\]

Then apply Exer. 13 of Ch. IV, §5. If B is \(\mu\) -measurable, \(\mu^{*}(\mathbf{A}) < +\infty\) and \((\mathbf{A}_{n})\) is a decreasing sequence of integrable subsets of G containing A such that \(\inf (\mu (\mathbf{A}_n)) = \mu^{*}(\mathbf{A})\) , show that

\[
\mu^ {*} (s \mathrm{A} \cap \mathrm{B}) = \inf \left(\mu^ {*} (s \mathrm{A} _ {n} \cap \mathrm{B})\right)
\]

(Ch. V, 1st edn., §2, No.~2, Lemma 1); \(^{4}\) note on the other hand that \(\mu(A_{n}-A_{n+1})\) tends to 0 with \(1/n\) .)

\begin{enumerate}
\def\labelenumi{\alph{enumi})}
\setcounter{enumi}{1}
\item
  If \(\mathbf{A}^{-1}\) is \(\mu\) -integrable and \(\mu^{*}(\mathbf{B}) < +\infty\) , then the function \(f\) is uniformly continuous for both the right and left uniform structures of \(\mathbf{G}\) ; moreover, one then has \(\int_{\mathbf{G}} f(s) d\mu(s) = \mu(\mathbf{A}^{-1})\mu^{*}(\mathbf{B})\) . (Reduce to the case that \(\mathbf{B}\) is integrable; observe then that \(\mu(s\mathbf{A} \cap \mathbf{B}) = \mu(\mathbf{A} \cap s^{-1}\mathbf{B})\) , that \(\varphi_{s\mathbf{A} \cap \mathbf{B}} = \varphi_{s\mathbf{A}}\varphi_{\mathbf{B}}\) and that \(\varphi_{s\mathbf{A}}(t) = \varphi_{t\mathbf{A}^{-1}}(s)\) .)
\item
  Deduce from \(a\) ) that in the two cases considered, the interiors of AB and BA are not empty if A and B are not negligible. (Cf. Ch. VIII, §4, No.~6, Prop. 17.)
\item
  In the group \(\mathbf{G} = \mathbf{SL}_{2}(\mathbf{R})\) , give an example of a compact set A and a \(\mu\) -measurable set B such that \(f(s) = \mu(sA \cap B)\) is not uniformly continuous for the left uniform structure. (Observe that there exist a sequence \((t_{n})\) of elements of G tending to e and a sequence \((s_{n})\) of elements of G such that the sequence \((s_{n}^{-1}t_{n}s_{n})\) tends to the point at infinity.)
\item
  Give an example of two nonmeasurable sets A, B in a locally compact group G, of finite outer measure and such that the function \(s \mapsto \mu^{*}(sA \cap B)\) is not continuous (cf.~Ch. IV, §4, Exer. 8).
\end{enumerate}

\begin{enumerate}
\def\labelenumi{\arabic{enumi})}
\setcounter{enumi}{20}
\tightlist
\item
  \begin{enumerate}
  \def\labelenumii{\alph{enumii})}
  \tightlist
  \item
    Let G be a locally compact group, \(\mu\) a left Haar measure on G. Show that if S is a stable subset of G such that \(\mu_{*}(S)>0\) , then the interior of S is nonempty (cf.~Exer. 20). In particular, if a subgroup H of G has nonzero inner measure, then H is an open subgroup of G.
  \end{enumerate}
\end{enumerate}

\begin{enumerate}
\def\labelenumi{\alph{enumi})}
\setcounter{enumi}{1}
\item
  If G is compact, then every stable subset S of G such that \(\mu_{*}(S) > 0\) is a compact open subgroup of G (observe that the interior of S is stable, and make use of Exer. 17 b)).
\item
  Assume that G is compact and abelian, written additively, and that there exists in G an element a of infinite order. Show that there exists a stable subset S of G such that \(a \in S\) , \(-a \notin S\) and \(G = S \cup (-S)\) (use Zorn's lemma); deduce from b) that S is not measurable and that \(\mu_{*}(S) = 0\) .
\end{enumerate}

\begin{enumerate}
\def\labelenumi{\arabic{enumi})}
\setcounter{enumi}{21}
\tightlist
\item
  Let G be a locally compact group, \(\mu\) a left Haar measure on G, and A an integrable subset of G such that \(\mu(A) > 0\) . Show that the set H(A) of \(s \in G\) such that \(\mu(A) = \mu(A \cap sA)\) is a compact group. (Observe that H(A) is closed in G, with the help of Exer. 20. To see that H(A) is compact, consider a compact subset B of A such that \(\mu(B) > \mu(A)/2\) and prove that \(H(A) \subset BB^{-1}\) .)
\end{enumerate}

¶ 23) Let G be a locally compact abelian group, written additively, μ a Haar measure on G, and A,B two integrable subsets of G.

\begin{enumerate}
\def\labelenumi{\alph{enumi})}
\tightlist
\item
  For every \(s \in \mathbf{G}\) , let
\end{enumerate}

\[
\mathrm{A} ^ {\prime} = \sigma_ {s} (\mathrm{A}, \mathrm{B}) = \mathrm{A} \cup (\mathrm{B} + s), \quad \mathrm{B} ^ {\prime} = \tau_ {s} (\mathrm{A}, \mathrm{B}) = (\mathrm{A} - s) \cap \mathrm{B}.
\]

Show that \(\mu (\mathrm{A}^{\prime}) + \mu (\mathrm{B}^{\prime}) = \mu (\mathrm{A}) + \mu (\mathrm{B})\) and \(\mathrm{A}' + \mathrm{B}'\subset \mathrm{A} + \mathrm{B}\)

\begin{enumerate}
\def\labelenumi{\alph{enumi})}
\setcounter{enumi}{1}
\tightlist
\item
  Assume that 0 belongs to \(\mathbf{A} \cap \mathbf{B}\) . A pair \((\mathbf{A}', \mathbf{B}')\) of integrable subsets of \(\mathbf{G}\) is said to be derived from (A, B) if there exist a sequence \((s_k)_{1 \leqslant k \leqslant n}\) of elements of \(\mathbf{G}\) and two sequences \((\mathbf{A}_k)_{0 \leqslant k \leqslant n}\) , \((\mathbf{B}_k)_{0 \leqslant k \leqslant n}\) of subsets of \(\mathbf{G}\) such that \(\mathbf{A}_0 = \mathbf{A}\) , \(\mathbf{B}_0 = \mathbf{B}\) , \(\mathbf{A}_k = \sigma_{s_k}(\mathbf{A}_{k-1}, \mathbf{B}_{k-1})\) , \(\mathbf{B}_k = \tau_{s_k}(\mathbf{A}_{k-1}, \mathbf{B}_{k-1})\) for \(1 \leqslant k \leqslant n\) , \(s_k \in \mathbf{A}_{k-1}\) for \(1 \leqslant k \leqslant n\) , and \(\mathbf{A}' = \mathbf{A}_n\) , \(\mathbf{B}' = \mathbf{B}_n\) . Show that there exists a sequence of pairs \((\mathrm{E}_n, \mathrm{F}_n)\) such that \(\mathbf{E}_0 = \mathbf{A}\) , \(\mathbf{F}_0 = \mathbf{B}\) , \((\mathrm{E}_{n+1}, \mathrm{F}_{n+1})\) is derived from \((\mathrm{E}_n, \mathrm{F}_n)\) , and \(\mu((\mathrm{E}_n - s) \cap \mathrm{F}_n) \geqslant \mu(\mathrm{F}_{n+1}) - 2^{-n}\) for every \(n\) and every \(s \in \mathrm{E}_n\) . Set \(\mathbf{E}_{\infty} = \bigcup_{n} \mathbf{E}_n\) , \(\mathbf{F}_{\infty} = \bigcap_{n} \mathbf{F}_n\) . Show that for every \(s \in \mathrm{E}_{\infty}\) ,
\end{enumerate}

\[
\mu \big ((\mathrm{E} _ {\infty} - s) \cap \mathrm{F} _ {\infty} \big) = \mu (\mathrm{F} _ {\infty}).
\]

\begin{enumerate}
\def\labelenumi{\alph{enumi})}
\setcounter{enumi}{2}
\tightlist
\item
  Assume that \(\mu (\mathrm{F}_{\infty}) > 0\) . Show that the function
\end{enumerate}

\[
f (s) = \mu \big ((\mathrm{E} _ {\infty} - s) \cap \mathrm{F} _ {\infty} \big)
\]

can only take on the values 0 and \(\mu(\mathrm{F}_{\infty})\) , and that the set C of \(s \in G\) such that \(f(s) = \mu(\mathrm{F}_{\infty})\) is open and closed, is such that \(\mu(C) = \mu(\mathrm{E}_{\infty})\) , and is the closure of \(\mathrm{E}_{\infty}\) (use Exer. 20 a) and b)). On the other hand let D be the set of \(s \in \mathrm{F}_{\infty}\) such that the intersection of \(\mathrm{F}_{\infty}\) with every neighborhood of \(s\) has measure \(>0\) . Show that \(\mu(D) = \mu(\mathrm{F}_{\infty})\) and that

\[
\mathbf {E} _ {\infty} + \mathbf {D} \subset \mathbf {C};
\]

from this, deduce that D is contained in the subgroup H(C) defined in Exer. 22, and that H(C) is compact and open in G. Finally, show that C + H(C) = C, that \(\mu(C) \geqslant \mu(A) + \mu(B) - \mu(H(C))\) , and that C ⊂ A + B (consider the measure of E∞ ∩ (c - F∞) for every c ∈ C).

\begin{enumerate}
\def\labelenumi{\alph{enumi})}
\setcounter{enumi}{3}
\tightlist
\item
  Deduce from c) that, for two integrable subsets A, B of G: either \(\mu_{*}(A+B)\geqslant\mu(A)+\mu(B)\) ; or else there exists a compact open subgroup H of G such that \(A+B\) contains a coset of H, in which case \(\mu_{*}(A+B)\geqslant\mu(A)+\mu(B)-\mu(H)\) . The case that G is connected.
\end{enumerate}

¶ 24 a) In R, let A (resp. B) be the set of numbers x whose proper dyadic expansion \(x = x_0 + \sum_{i=1}^{\infty} x_i 2^{-i}\) ( \(x_0\) an integer, \(x_i = 0\) or \(x_i = 1\) for \(i \geqslant 1\) ) is such that \(x_{i}=0\) for i even and \textgreater0 (resp. i odd). Show that, for Lebesgue measure, A and B have measure zero but \(A+B=R\) .

\begin{enumerate}
\def\labelenumi{\alph{enumi})}
\setcounter{enumi}{1}
\item
  Deduce from a) that there exists a basis H of R (over Q) contained in A ∪ B, hence of measure zero. The set P₁ of numbers rh, where r ∈ Q and h ∈ H, is also of measure zero.
\item
  Denote by \(P_{n}\) the set of real numbers at most n of whose coordinates relative to the basis H are nonzero. Show that if \(P_{n}\) is negligible and \(P_{n+1}\) is measurable, then \(P_{n+1}\) is negligible. (Let \(h_{0} \in H\) ; show first that the set S of \(x \in P_{n+1}\) whose coordinate relative to \(h_{0}\) is \(\neq 0\) , is negligible. Using Exer. 20, show that if \(P_{n+1}\) were not negligible, there would exist two distinct points \(x', x''\) of \(P_{n+1} \cap CS\) such that \((x' - x'') / h_{0}\) is rational, and deduce from this a contradiction.)
\item
  Deduce from a) and b) that there exist in R two negligible sets C, D such that C + D is not measurable.
\end{enumerate}

¶ 25) a) Let f be a positive numerical function defined on R, integrable (for the Lebesgue measure μ on R), bounded, and with compact support. Let γ = sup f(t).

For every \(w \in \mathbf{R}\) , denote by \(\mathrm{U}_f(w)\) the set of \(t \in \mathbf{R}\) such that \(f(t) \geqslant w\) ; set \(\nu_f(w) = \mu^*(\mathrm{U}_f(w))\) . Show that for every \(\alpha > 1\) ,

\[
\int_ {- \infty} ^ {+ \infty} f ^ {\alpha} (t) d t = \int_ {0} ^ {\gamma} \nu_ {f} (w) \alpha w ^ {\alpha - 1} d w.
\]

\begin{enumerate}
\def\labelenumi{\alph{enumi})}
\setcounter{enumi}{1}
\tightlist
\item
  Let g be a second numerical function satisfying the same conditions as f, and set \(\delta = \sup_{t \in \mathbf{R}} g(t)\) . Let h be the function defined on \(R^{2}\) by \(h(u,v) = f(u) + g(v)\) if \(f(u)g(v) \neq 0\) , and \(h(u,v) = 0\) otherwise; finally, set \(k(t) = \sup_{u+v=t} h(u,v)\) , so that k is positive, integrable, bounded, and has compact support. Show that for every \(\alpha > 1\) ,
\end{enumerate}

\[
\int_ {- \infty} ^ {+ \infty} k ^ {\alpha} (t) d t \geqslant (\gamma + \delta) ^ {\alpha} \left(\frac {1}{\gamma^ {\alpha}} \int_ {- \infty} ^ {+ \infty} f ^ {\alpha} (t) d t + \frac {1}{\delta^ {\alpha}} \int_ {- \infty} ^ {+ \infty} g ^ {\alpha} (t) d t\right).
\]

(Observe that for \(0 < w < 1\) , one has \(\mathrm{U}_k(\gamma w + \delta w) \supset \mathrm{U}_f(\gamma w) + \mathrm{U}_g(\delta w)\) , and use \(a\) and Exer. 23 d).)

\begin{enumerate}
\def\labelenumi{\alph{enumi})}
\setcounter{enumi}{2}
\tightlist
\item
  Let \(\mu_{n}\) be Lebesgue measure on \(\mathbf{R}^{n}\) , A and B two integrable subsets of \(\mathbf{R}^{n}\) . Show that (Brunn-Minkowski inequality)
\end{enumerate}

\[
\left(\left(\mu_ {n}\right) _ {*} (\mathrm{A} + \mathrm{B})\right) ^ {1 / n} \geqslant \left(\mu_ {n} (\mathrm{A})\right) ^ {1 / n} + \left(\mu_ {n} (\mathrm{B})\right) ^ {1 / n}.
\]

(Reduce to the case that A and B are compact. Then argue by induction on n using Exer. 23 d), the Lebesgue--Fubini theorem, the inequality proved in b), as well as Hölder's inequality.)

¶ 26) Let G be a locally compact group, μ a left Haar measure on G, k an integer ≥1, and A an integrable set. Show that for every ε \textgreater{} 0, there exists a neighborhood U of e in G having the following property: for every finite subset S of k elements in U, the set of s ∈ G such that sS ⊂ A has measure ≥(1 - ε)μ(A). (Reduce to the case that A is compact and take U to be such that

\[
\mu (\mathrm{AU}) \leqslant \left(1 - \frac {\varepsilon}{k - 1}\right) \mu (\mathrm{A}).
\]

For every finite subset H of G, set \(p(\mathrm{H}) = \operatorname{Card}(\mathrm{H})\) and \(q(\mathrm{H}) = 1\) if \(\mathrm{H} \neq \emptyset\) , \(q(\mathrm{H}) = 0\) if \(\mathrm{H} = \emptyset\) ; by evaluating the integrals

\[
\int_ {\mathbf {G}} p (\mathrm{A} \cap s \mathrm{S}) d s \quad \text { and } \quad \int_ {\mathbf {G}} q (\mathrm{A} \cap s \mathrm{S}) d s  ,
\]

show that if, for \(h \leqslant k\) , \(\mathbf{M}_h\) denotes the set of \(s \in G\) for which \(\operatorname{Card}(\mathbf{A} \cap s\mathbf{S}) = h\) , then

\[
\sum_ {h = 1} ^ {k} h \mu (\mathrm{M} _ {h}) = k \cdot \mu (\mathrm{A}) \quad \text { and } \quad \sum_ {h = 1} ^ {k} \mu (\mathrm{M} _ {h}) \leqslant \mu (\mathrm{AU}).
\]

From this, conclude that \(\sum_{h=1}^{k-1} h\mu(M_h) \leqslant k\varepsilon \cdot \mu(A)\) .

\begin{enumerate}
\def\labelenumi{\arabic{enumi})}
\setcounter{enumi}{26}
\tightlist
\item
  Let \(\mathbf{G}\) be a group operating on a set \(\mathbf{X}\) . A subset \(\mathbf{P}\) (resp. \(\mathbf{C}\) ) of \(\mathbf{X}\) is said to be a G-filling (resp. G-covering) if, for every \(s \neq e\) in \(\mathbf{G}\) , one has \(s\mathbf{P} \cap \mathbf{P} = \emptyset\) (resp. if \(\mathbf{X} = \bigcup_{s \in \mathbf{G}} s\mathbf{C}\) ). One calls G-paving a subset \(\mathbf{P}\) that is both a G-filling and a G-covering.
\end{enumerate}

\begin{enumerate}
\def\labelenumi{\alph{enumi})}
\item
  Assume that X is locally compact, that G is countable, operates continuously in X, and that there exists a nonzero positive measure \(\mu\) on X invariant under G. Let P and C be a G-filling and a G-covering that are \(\mu\) -integrable. Show that \(\mu(C) \geqslant \mu(P)\) . (Observe that \(\mu(C) \geqslant \sum_{s \in G} \mu(C \cap sP)\) .)
\item
  Assume in addition that there exists on X a uniform structure compatible with the topology of X and admitting a fundamental system S of open entourages invariant under G. On the other hand, one denotes by \(\Delta(G)\) the infimum of the numbers \(\mu(C)\) over all the integrable G-coverings C of X. Let V be an entourage belonging to S, and let \(a \in X\) be such that \(\mu(V(a)) > \Delta(G)\) ; show that there exists an \(s \in G\) such that \(s \neq e\) and \((a, sa) \in \overline{V} \circ V\) .
\item
  Assume that X is a locally compact group, \(\mu\) a left Haar measure, and G a countable subgroup of X, operating by left translation. With the same definition of \(\Delta(G)\) as in \(b\) ), show that if A is an integrable subset of X such that \(\mu(A) > \Delta(G)\) , then there exists \(s \in G \cap AA^{-1}\) such that \(s \neq e\) . Special case that \(X = R^n\) and G is a discrete subgroup of rank \(n\) in \(\mathbf{R}^n: \Delta(G)\) is then equal to the absolute value of the determinant of a basis of G over Z with respect to the canonical basis of \(\mathbf{R}^n\) ; if A is a symmetric closed convex set with nonempty interior \(^5\) in \(\mathbf{R}^n\) such that \(\mu(A) \geqslant 2^n \Delta(G)\) , show that there exists a point of \(A \cap G\) distinct from 0 (Minkowski's theorem).
\item
  With hypotheses as in a), let \(f\) be a function \(\geqslant 0\) and \(\mu\) -integrable on \(X\) . Show that there exist two points \(a, b\) of \(X\) such that \(\mu(C) \sum_{s \in G} f(sa) \geqslant \int_X f(x) d\mu(x)\) and \(\mu (\mathrm{P})\sum_{s\in \mathbb{G}}f(sb)\leqslant \int_{\mathbb{X}}f(x)d\mu (x).\) (Observe that if \(g\) is an integrable function \(\geqslant 0\) , E a
\end{enumerate}

\(\mu\) -integrable set in \(\mathbf{X}\) , then there exists a \(c \in \mathbf{E}\) such that \(\int_{\mathbf{E}} g(x) d\mu(x) \leqslant g(c)\mu(\mathbf{E})\) , and a \(c' \in \mathbf{E}\) such that \(\int_{\mathbf{E}} g(x) d\mu(x) \geqslant g(c')\mu(\mathbf{E})\) .

¶ 28) With hypotheses as in Exer. 27 a), assume in addition that there exists a μ-integrable G-paving F.

\begin{enumerate}
\def\labelenumi{\alph{enumi})}
\tightlist
\item
  For every \(\mu\) -integrable numerical function \(f \geqslant 0\) defined on X, set \(\widetilde{f}(x) = \sum_{s \in G} f(sx)\) , so that \(\int_{F} \widetilde{f}(x) d\mu(x) = \int_{X} f(x) d\mu(x)\) (§2, No.~10, Prop. 15). Let \((f_i)_{1 \leqslant i \leqslant n}\) be a family of numerical functions \(\geqslant 0\) , \(\mu\) -integrable on X: for \(i \neq j\) set
\end{enumerate}

\[
m _ {i j} = \int_ {\mathbf {F}} \widetilde {f} _ {i} (x) \widetilde {f} _ {j} (x) d \mu (x),
\]

and let \(c_{i} = \sup_{x\in X}\widetilde{f}_{i}(x)\) . Show that there exists at least one pair of distinct indices \((i,j)\) such that

\[
m _ {i j} \geqslant \frac {1}{n (n - 1) \mu (\mathrm{F})} \left(\left(\sum_ {i = 1} ^ {n} \int_ {\mathrm{X}} f _ {i} (x) d \mu (x)\right) ^ {2} - \mu (\mathrm{F}) \sum_ {i = 1} ^ {n} c _ {i} \int_ {\mathrm{X}} f _ {i} (x) d x\right).
\]

(Bound the sum \(\sum_{i\neq j}m_{ij}\) from below, using the Cauchy-Schwarz inequality.) From this, deduce that if \((\mathrm{A}_i)_{1\leqslant i\leqslant n}\) is a finite family of \(\mu\) -integrable G-fillings, and if

\[
\sum_ {i = 1} ^ {n} \mu (\mathrm{A} _ {i}) > \mu (\mathrm{F}),
\]

then there exist a pair \((i,j)\) of distinct indices and an \(s\in G\) such that \(\mu (\mathrm{A}_is\cap \mathrm{A}_j) > 0\) b) If \(\mathbf{G}_0\) is a subgroup of \(G\) of finite index \((G:\mathbf{G}_0) = h\) , and if \((s_1,\ldots ,s_h)\) is a system of representatives of the right cosets of \(\mathbf{G}_0\) in \(G\) , then \(\mathrm{F_0} = \bigcup_{1\leqslant i\leqslant h}s_i\mathrm{F}\) is a \(G_{0}\)-paving.

\begin{enumerate}
\def\labelenumi{\alph{enumi})}
\setcounter{enumi}{2}
\tightlist
\item
  In \(\mathbf{X} = \mathbf{R}^n\) , let \(\mathbf{A}\) be a symmetric closed convex set with nonempty interior; let \(u_i : (x_j) \mapsto \sum_{j=1}^{n} c_{ij} x_j\) be \(m\) linear forms on \(\mathbf{X}\) , with integer coefficients \(c_{ij} (m < n)\) .
\end{enumerate}

Show that for every integer \(p \geqslant 1\) and every number \(r \geqslant 0\) such that \(\mu(A)r^n \geqslant 2^n p^m\) , there exists a point \(x \in rA \cap Z^n\) distinct from 0 and such that \(u_i(x) \equiv 0 \pmod{p}\) for \(1 \leqslant i \leqslant m\) (apply Minkowski's theorem (Exer. 27 c)) to the subgroup \(G_0\) of \(Z^n\) formed by the \(z \in Z^n\) such that \(u_i(z) \equiv 0 \pmod{p}\) for \(1 \leqslant i \leqslant m\) , and make use of \(b\) ). Special case that \(n = 2\) , \(m = 1\) and A is defined by \(|x_1| \leqslant 1\) , \(|x_2| \leqslant 1\) (Thue's theorem).

¶ 29) a) Let p be a prime number; there exist two integers a, b such that \(a^{2}+b^{2}+1 \equiv 0 \pmod{p}\) (Alg., Ch. V, 1st edn., §11, No.~5, Cor. of Th. 3). Show that there exist integers \(x_{1}, x_{2}, x_{3}, x_{4}\) not all zero such that \(ax_{1} + bx_{2} \equiv x_{3} \pmod{p}\) , \(bx_{1} - ax_{2} \equiv x_{4} \pmod{p}\) and

\[
y = x _ {1} ^ {2} + x _ {2} ^ {2} + x _ {3} ^ {2} + x _ {4} ^ {2} \leqslant \sqrt {2} p
\]

(same method as in Exer. 28 c)). Show that \(y\) is divisible by \(p\) , and deduce therefrom that \(y = p\) .

\begin{enumerate}
\def\labelenumi{\alph{enumi})}
\setcounter{enumi}{1}
\tightlist
\item
  Deduce from a) that every integer \(n \geqslant 0\) is the sum of four squares (Lagrange's theorem; make use of Alg., Ch. IV, 1st edn., §2, Exer. 11).
\end{enumerate}

\begin{enumerate}
\def\labelenumi{\arabic{enumi})}
\setcounter{enumi}{29}
\tightlist
\item
  Let G be a locally compact group, \(\mu\) a left Haar measure on G, and \(\nu\) a bounded measure on G. Assume that the mapping \(s \mapsto \gamma(s)\nu\) of G into the Banach space \(\mathcal{M}^{1}(G)\) is continuous. Show that the measure \(\nu\) has base \(\mu\) . (Let A be a \(\mu\) -negligible compact subset of G. Arguing as for Prop. 11, show that \(\nu(xA)=0\) for x running over a dense subset of G. From this, deduce that \(\nu(A)=0\) .)
\end{enumerate}

Conversely, if \(\nu\) has base \(\mu\) , then the mapping \(s \mapsto \pmb{\gamma}(s)\nu\) of G into \(\mathcal{M}^1(\mathrm{G})\) is continuous: cf.~Ch. VIII, §2, No.~5, Prop. 8.
